\clearpage
\onecolumn
\thispagestyle{empty}
\vspace*{\fill}
\begin{center}
	{\color{radblue}\rule{\linewidth}{2pt}}\par
	\vspace{6mm}
	{\Huge\bfseries\color{radblue} Pelvis}\par
	\vspace{6mm}
	{\color{radblue}\rule{\linewidth}{2pt}}
\end{center}
\vspace*{\fill}
\clearpage
\twocolumn
\pagestyle{contenido}


%%%%%%%%%%%%%%%%%%%%%%%%%%%%%%%%%%%%%%%%%%%%%%%%%%%%%%%%%%%
%  SECCIÓN: Pelvis — Introducción Anatómica
%  Archivo: secciones/pelvis.tex
%  Sin preámbulo ni \begin{document} — lo maneja main.tex
	%%%%%%%%%%%%%%%%%%%%%%%%%%%%%%%%%%%%%%%%%%%%%%%%%%%%%%%%%%%
	
	\section{Pelvis}
	
	\subsection{Anatomía de la Región}
	
	\subsubsection*{Composición Ósea}
	
	La pelvis está formada por \textbf{cuatro huesos}: dos huesos coxales, el sacro y el cóccix.
	Cada hueso coxal resulta de la fusión de tres piezas: \textbf{ilion}, \textbf{isquion} y \textbf{pubis}.
	En toda radiografía de pelvis se incluye obligatoriamente la \textbf{porción proximal de ambos fémures},
	con un rango que abarca desde las crestas ilíacas hasta por debajo del trocánter menor.
	
	\subsubsection*{Fémur Proximal}
	
	\paragraph{Componentes}
	La extremidad proximal del fémur comprende: \textbf{cabeza femoral} (con la fóvea capitata, depresión
	para la inserción del ligamento redondo), \textbf{cuello femoral}, \textbf{trocánter mayor}
	(superior y lateral) y \textbf{trocánter menor} (posterior y medial).
	
	\paragraph{Angulaciones del cuello femoral}
	El cuello femoral se proyecta en dirección medial, superior y \textbf{anterior 15--20°} respecto
	a la diáfisis en el plano axial. En el plano longitudinal forma un ángulo de \textbf{120--130°}
	con la diáfisis. Estos valores varían según edad, sexo y estatura, y son el fundamento de la
	rotación interna de extremidades durante el posicionamiento.
	
	\subsubsection*{Articulaciones}
	
	\textbf{Sacroilíacas:} unen el sacro con los huesos ilíacos.
	
	\textbf{Sínfisis púbica:} unión fibrocartilaginosa de los dos pubis en la línea media.
	
	\textbf{Coxofemorales (caderas):} articulación entre el acetábulo y la cabeza femoral.
	
	\subsubsection*{Diferencias entre Pelvis Femenina y Masculina}
	
	La \textbf{pelvis femenina} es más ligera, ancha y menos profunda; presenta una entrada
	mayor y más ovalada, y un sacro más ancho, curvado posteriormente y más plano.
	La \textbf{pelvis masculina} es más pesada, estrecha y profunda, con una entrada menor.
	Estas diferencias morfológicas condicionan la técnica radiográfica y la interpretación
	de la simetría en la imagen.
	
	\subsubsection*{Reperes Anatómicos}
	
	\paragraph{Puntos de referencia palpables}
	Crestas ilíacas, espinas ilíacas anterosuperiores (EIAS), sínfisis púbica, trocánteres
	mayores, tuberosidades isquiáticas y vértice del cóccix.
	
	\paragraph{Localización de estructuras}
	\textbf{Trocánter mayor y sínfisis púbica} se encuentran en el mismo plano horizontal;
	esto permite localizar la sínfisis palpando el trocánter mayor y trazando una línea
	transversal. El \textbf{acetábulo} se localiza en el punto medio entre la EIAS y la
	sínfisis púbica; una línea perpendicular a esa referencia coincide con el eje
	longitudinal del cuello femoral.
	
	
	\begin{figure}[h]
		\centering
		\includegraphics[width=0.7\linewidth]{imagenes/pelvis1}
		\caption{}
		\label{fig:pelvis1}
	\end{figure}
	
	\begin{figure}[h]
		\centering
		\includegraphics[width=0.7\linewidth]{imagenes/pelvis2}
		\caption{}
		\label{fig:pelvis2}
	\end{figure}
	
	\begin{figure}
		\centering
		\includegraphics[width=0.7\linewidth]{imagenes/pelvis3}
		\caption{}
		\label{fig:pelvis3}
	\end{figure}
	
	
	\clearpage
	
	%----------------------------------------------------------
	%  Continúa con los enfoques desde \subsection{}
	%----------------------------------------------------------

%%%%%%%%%%%%%%%%%%%%%%%%%%%%%%%%%%%%%%%%%%%%%%%%%%%%%%%%%%%
%  SECCIÓN: Pelvis
%  Archivo: secciones/pelvis.tex
%  Sin preámbulo ni \begin{document} — lo maneja main.tex
	%%%%%%%%%%%%%%%%%%%%%%%%%%%%%%%%%%%%%%%%%%%%%%%%%%%%%%%%%%%
	
	\section{Pelvis}
	
	\subsection{Pelvis AP}
	
	\subsubsection*{1. Identificación}
	
	\textbf{Tipo:} Proyección básica. Forma par con la proyección oblicua (Cleaves modificado).
	
	\textbf{Indicaciones clínicas principales:}
	\begin{itemize}
		\item Fracturas y luxaciones articulares de la cintura pélvica y caderas
		\item Artropatías degenerativas y artrosis (coxartrosis)
		\item Evaluación de prótesis de cadera y dolor referido a la región pélvica
		\item Desniveles pélvicos, asimetrías de miembros inferiores e inestabilidad pélvica en bipedestación
	\end{itemize}
	
	\subsubsection*{2. Técnica Base}
	
	\textbf{Receptor de imagen:} 35 $\times$ 43 cm. Orientación \textbf{transversal} en paciente estándar.
	
	\textbf{DFRI:} 100 cm. Distancia estándar que minimiza la divergencia del haz y la magnificación de las estructuras pélvicas profundas.
	
	\textbf{Técnica:} Intermedia, equilibrio entre escala larga de grises para visualizar el trabeculado óseo del cuello femoral y alas ilíacas, y suficiente penetración para atravesar la densidad del sacro y la articulación coxofemoral. Una radiografía radiopaca indica falta de kVp; una radiografía radiolúcida indica falta de mAs.
	
	\textbf{Factores técnicos:} kVp \textbf{[según biotipo]} --- mAs \textbf{[según biotipo]}. Medir el espesor en la región de mayor densidad (acetábulo y sacro).
	
	\textbf{Bucky:} Bucky mural o de mesa con grilla, obligatorio dado el espesor de la región. Técnica homogenizada en pacientes pediátricos.
	
	\subsubsection*{3. Posicionamiento}
	
	\paragraph{Paciente}
	Decúbito supino o bipedestación si se solicita expresamente. Manos apoyadas sobre el tórax o brazos en posición cómoda separados del cuerpo. El plano medio sagital se orienta perpendicular al receptor y centrado en la línea media del Bucky.
	
	\paragraph{Región de interés}
	EIAS equidistantes a cada lado de la mesa para confirmar ausencia de rotación pélvica. De no estar contraindicada la movilización (fractura), rotar internamente pies y extremidades inferiores 15--20° para colocar los cuellos femorales paralelos al receptor. Los talones se separan aproximadamente 20--24 cm. Inmovilizar con sacos de arena si es necesario.
	
	\begin{figure}[h]
		\centering
		\includegraphics[width=0.7\linewidth]{imagenes/pelvis5}
		\caption{}
		\label{fig:pelvis5}
	\end{figure}
	
	
	\paragraph{Rayo Central}
	Perpendicular al plano del receptor. Punto de entrada: \textbf{punto medio entre el nivel de las EIAS y la sínfisis pubiana, aproximadamente 5 cm por debajo del nivel de las EIAS}.
	
	\paragraph{Receptor de imagen}
	Centrado en la línea media del paciente a la altura del punto de entrada del RC. Debe incluir toda la cintura pelviana desde las crestas ilíacas en su totalidad hasta el tercio superior de ambos fémures, y lateralmente las partes blandas cutáneas.
	
	\paragraph{Respiración}
	Suspendida en apnea al momento de la exposición para eliminar el movimiento involuntario de las estructuras abdominopélvicas y garantizar nitidez de la imagen.
	
	\subsubsection*{4. Anatomía Visible}
	
	Toda la pelvis ósea con sacro y cóccix, articulaciones sacroilíacas bilaterales, acetábulos, sínfisis púbica y la porción proximal de ambos fémures incluyendo cabeza, cuello, trocánter mayor y trocánter menor. Las articulaciones coxofemorales, sacroilíacas y la sínfisis púbica deben estar abiertas y evaluables. La simetría pélvica debe permitir comparar alas ilíacas y agujeros obturadores de ambos lados. El cuello femoral debe visualizarse en su extensión completa y paralelo al receptor; el trocánter mayor debe aparecer de perfil y el trocánter menor no debe verse o solo en su extremo medial.
	

	\subsubsection*{5. Criterios de Evaluación}
	
	\textbf{Posicionamiento correcto:} Toda la pelvis junto con la parte proximal de los fémures incluida en el campo. Cuellos femorales desplegados en toda su extensión sin superposición. Trocánteres mayores de perfil con tamaño y forma simétricos. Sacro y cóccix alineados con la sínfisis del pubis. Trocánter mayor al mismo nivel que la sínfisis del pubis.
	
	\textbf{Ausencia de rotación:} Agujeros obturadores simétricos y de igual tamaño. Alas ilíacas simétricas. Espinas isquiáticas visibles con igual tamaño. Ambos ilíacos y ambos trocánteres mayores equidistantes de los márgenes del receptor. Parte inferior de la columna vertebral centrada en la imagen.
	
	\textbf{Penetración y contraste:} Trama ósea trabecular visible en cuello femoral y alas ilíacas. Partes blandas adyacentes y región vesical diferenciables. Densidad del sacro y articulación coxofemoral correctamente penetrada sin sobreexposición de las alas ilíacas.
	
	\subsubsection*{6. Puntos Críticos}
	
	\textbf{Rotación pélvica:} Un agujero obturador más cerrado indica rotación del paciente hacia ese mismo lado; el ala ilíaca ipsilateral se verá más estrecha y de perfil mientras la contralateral aparece más ancha. Corregir reposicionando las EIAS a igual distancia de la mesa antes de la exposición.
	
	\textbf{Omisión de la rotación interna de extremidades:} Sin la rotación interna de 15--20°, el cuello femoral se acorta por su angulación anterior de 15--20° respecto a la diáfisis y el trocánter menor se proyecta en su totalidad sobre el borde medial del fémur. Con rotación externa el cuello se acorta y el trocánter menor es claramente visible; con la rotación interna correcta el cuello se despliega y el trocánter menor apenas se insinúa o desaparece.
	
	\textbf{Adaptaciones:} En sospecha de fractura no realizar rotación interna de extremidades. En bipedestación, indicada para desniveles pélvicos, inestabilidad o evaluación postural; el reparo del trocánter mayor se utiliza pidiendo al paciente que junte las puntas de los pies y separe los talones. En pacientes con prótesis de cadera, consultar si existe dolor agudo o riesgo de luxación antes de realizar la rotación interna; la prótesis reproduce la morfología del hueso y los cuellos quedan igualmente desplegados al rotar.
	
	\textbf{Trucos:} Para localizar la sínfisis púbica, palpar el trocánter mayor y trazar una línea horizontal: ambas estructuras se encuentran en el mismo plano. El acetábulo se localiza en el punto medio entre la EIAS y la sínfisis púbica. En la bipedestación, pedir al paciente que apoye igual peso en ambos miembros para evitar desniveles posicionales que simulen patología.
	
	\clearpage
	\onecolumn
	
	\begin{figure}
		\centering
		\includegraphics[width=0.7\linewidth]{imagenes/pelvis4}
		\caption{}
		\label{fig:pelvis4}
	\end{figure}
	
		\begin{figure}
		\centering
		\includegraphics[width=0.7\linewidth]{imagenes/pelvis6}
		\caption{}
		\label{fig:pelvis6}
	\end{figure}
		
	\twocolumn
	\clearpage
	
	
	%----------------------------------------------------------
	%  Para agregar otro enfoque duplicá el bloque
	%  desde \subsection hasta \textbf{Trucos}
	%----------------------------------------------------------
	
%%%%%%%%%%%%%%%%%%%%%%%%%%%%%%%%%%%%%%%%%%%%%%%%%%%%%%%%%%%
%  SECCIÓN: Pelvis Pediátrico
%  Archivo: secciones/pelvis_pediatrico.tex
%  Sin preámbulo ni \begin{document} — lo maneja main.tex
	%%%%%%%%%%%%%%%%%%%%%%%%%%%%%%%%%%%%%%%%%%%%%%%%%%%%%%%%%%%
	
	\section{Pelvis}
	
	\subsection{Pelvis AP Pediátrico}
	
	\subsubsection*{1. Identificación}
	
	\textbf{Tipo:} Proyección básica. Forma par con el perfil pediátrico (Lowenstein / Von Rosen).
	
	\textbf{Indicaciones clínicas principales:}
	\begin{itemize}
		\item Displasia del desarrollo de cadera
		\item Dislocación congénita de cadera
		\item Control evolutivo
	\end{itemize}
	
	\subsubsection*{2. Técnica Base}
	
	\textbf{Receptor de imagen:} 24$\times$30 cm. Orientación \textbf{transversal} en paciente pediátrico.
	
	\textbf{DFRI:} 100 cm.
	
	\textbf{Técnica:} Homogeneizada.
	
	\textbf{Bucky:} Sin Bucky ni grilla. En pacientes con espesores menores a 15 cm se suprime la rejilla para evitar el factor de multiplicación de dosis, cumpliendo el principio ALARA y protegiendo los órganos radiosensibles en desarrollo.
	
	\subsubsection*{3. Posicionamiento}
	
	\paragraph{Paciente}
	Decúbito supino. Retirar los pañales antes del posicionamiento.
	
	\paragraph{Región de interés}
	Desde las crestas ilíacas hasta el fémur proximal, bilateral.
	
	\paragraph{Rayo Central}
	Perpendicular al plano del receptor. Punto de entrada: \textbf{1 dedo por encima del pubis}.
	
	\paragraph{Receptor de imagen}
	Centrado a nivel del RC. Bordes superiores: crestas ilíacas completas. Bordes inferiores: fémur proximal bilateral.
	
	\paragraph{Respiración}
	\textit{Similar a adulto}.
	
	\subsubsection*{4. Anatomía Visible}
	
	\textit{Similar a adulto}.
	
	\subsubsection*{5. Criterios de Evaluación}
	
	\textbf{Posicionamiento correcto:} Las crestas ilíacas deben incluirse completas para permitir la medición de los grados de Risser (escala 0--5, osificación apofisaria de la cresta ilíaca como indicador de madurez esquelética).
	
	\textbf{Ausencia de rotación:} \textit{[dato no suministrado]}.
	
	\textbf{Penetración y contraste:} \textit{[dato no suministrado]}.
	
	\subsubsection*{6. Puntos Críticos}
	
	\textbf{Corte de crestas ilíacas:} Impide la medición de los grados de Risser.
	
	\textbf{\textit{[Error frecuente 2]:}} \textit{[dato no suministrado]}.
	
	\textbf{Adaptaciones:} \textit{[dato no suministrado]}.
	
	\textbf{Trucos:} Estudio bilateral y comparativo. La ausencia de Bucky en espesores menores a 15 cm es una medida de radioprotección activa.
	
	\clearpage
	
	%----------------------------------------------------------
	%  Para agregar otro enfoque duplicá el bloque
	%  desde \subsection hasta \textbf{Trucos}
	%----------------------------------------------------------


%%%%%%%%%%%%%%%%%%%%%%%%%%%%%%%%%%%%%%%%%%%%%%%%%%%%%%%%%%%
%  SECCIÓN: Cadera
%  Archivo: secciones/cadera.tex
%  Sin preámbulo ni \begin{document} — lo maneja main.tex
	%%%%%%%%%%%%%%%%%%%%%%%%%%%%%%%%%%%%%%%%%%%%%%%%%%%%%%%%%%%
	
	\section{Pelvis}
	
	\subsection{Cadera AP}
	
	\subsubsection*{1. Identificación}
	
	\textbf{Tipo:} Proyección básica. Forma par con la proyección lateral (axiolateral o Lauenstein).
	
	\textbf{Indicaciones clínicas principales:}
	\begin{itemize}
		\item Control posoperatorio y seguimiento de patología de cadera
		\item Evaluación del acetábulo, cabeza y cuello femoral
		\item Verificación de la colocación y estado de dispositivos ortopédicos (prótesis, clavos, tornillos)
	\end{itemize}
	
	\subsubsection*{2. Técnica Base}
	
	\textbf{Receptor de imagen:} 24×30 cm. Orientación \textbf{longitudinal}.
	
	\textbf{DFRI:} 100 cm. Distancia estándar para proyecciones de miembro inferior.
	
	\textbf{Técnica:} Homogeneizada. Permite equilibrar las diferencias de densidad entre las estructuras óseas de la región coxofemoral.
	
	\textbf{Bucky:} Bucky de mesa con grilla.
	
	\subsubsection*{3. Posicionamiento}
	
	\paragraph{Paciente}
	Decúbito supino sobre la mesa radiográfica. El paciente debe retirar la ropa de la región pélvica; se puede dejar en ropa interior o colocar chaleco de protección. El plano medio sagital queda perpendicular a la mesa.
	
	\paragraph{Región de interés}
	Pelvis sin rotación: las EIAS deben ser equidistantes a la mesa. El miembro inferior del lado a estudiar se rota internamente entre 15° y 20°, llevando el pie en rotación interna, con el objetivo de colocar el cuello femoral paralelo al plano del receptor. Esta maniobra está contraindicada en fracturas de cadera u otras situaciones clínicas que lo impidan. Se puede colocar un apoyo bajo las rodillas para aliviar la tensión lumbar del paciente.
	
	\begin{figure}[h]
		\centering
		\includegraphics[width=0.7\linewidth]{imagenes/pelvis7}
		\caption{}
		\label{fig:pelvis7}
	\end{figure}
	
	
	\paragraph{Rayo Central}
	Perpendicular al plano del receptor. Punto de entrada: \textbf{línea media entre la EIAS y la sínfisis del pubis, desplazado hacia distal — coincidente con el cuello femoral — o a nivel del pliegue inguinal}.
	
	\paragraph{Receptor de imagen}
	Centrado en la articulación coxofemoral. La región a incluir comprende el tercio proximal del fémur (cabeza, cuello y trocánteres) y parte del ilion y huesos del pubis. La cresta ilíaca no necesita incluirse obligatoriamente.
	
	\paragraph{Respiración}
	Suspendida en espiración. Evita el movimiento del paciente durante la exposición y minimiza la borrosidad por motilidad.
	
	\subsubsection*{4. Anatomía Visible}
	
	Se visualiza la articulación coxofemoral con la cabeza femoral proyectada a través del acetábulo. El cuello femoral debe aparecer sin acortamiento, con su eje longitudinal desplegado. El trocánter mayor se observa de perfil en el borde lateral. El trocánter menor no se visualiza o aparece de forma mínima en el borde medial, como consecuencia de la rotación interna aplicada. Se incluye parte del ilion supraacetabular y los huesos del pubis. En pacientes con dispositivos ortopédicos, estos deben visualizarse en su totalidad.
	
	\subsubsection*{5. Criterios de Evaluación}
	
	\textbf{Posicionamiento correcto:} Cuello femoral visible en su longitud real sin acortamiento; trocánter mayor de perfil; trocánter menor no visible o mínimamente visible en borde medial; articulación coxofemoral centrada en la imagen. Si hay prótesis u otro dispositivo ortopédico, debe verse completo dentro del campo.
	
	\textbf{Ausencia de rotación:} Las EIAS deben haberse posicionado equidistantes a la mesa antes de la exposición; ausencia de oblicuidad pélvica en la imagen resultante.
	
	\textbf{Penetración y contraste:} Técnica homogeneizada que permita diferenciar las estructuras óseas de la cabeza, cuello y acetábulo con adecuada escala de grises; ni sobreexpuesto ni subexpuesto en las trabéculas del cuello femoral.
	
	\subsubsection*{6. Puntos Críticos}
	
	\textbf{Omisión de rotación interna del miembro:} El cuello femoral queda en oblicuo, proyectándose acortado y sin desplegar su eje longitudinal real; corregir rotando el pie 15--20° hacia adentro, salvo contraindicación.
	
	\textbf{Rotación pélvica:} Las EIAS no quedan simétricas, generando asimetría del acetábulo y del ala ilíaca; corregir verificando equidistancia de EIAS antes de exponer.
	
	\textbf{Adaptaciones:} En paciente traumático sobre camilla o tabla de trauma, no movilizar el miembro para rotación interna; levantar al paciente con sábana y con asistencia médica para colocar el receptor. Si la rotación interna está contraindicada (fractura, prótesis luxada), documentar la limitación y obtener la imagen en posición neutra.
	
	\textbf{Trucos:} Para localizar el punto de entrada del RC: palpar EIAS, luego el trocánter mayor (a la misma altura se encuentra la sínfisis del pubis); el punto medio entre EIAS y sínfisis desplazado hacia distal corresponde al cuello femoral. Alternativamente, incidir directamente sobre el pliegue inguinal como referencia rápida.
	
	\clearpage
	\onecolumn
	
	\begin{figure}
		\centering
		\includegraphics[width=0.7\linewidth]{imagenes/pelvis8}
		\caption{}
		\label{fig:pelvis8}
	\end{figure}
	
	
	
	\twocolumn
	\clearpage
	
	%----------------------------------------------------------
	%  Para agregar otro enfoque duplicá el bloque
	%  desde \subsection hasta \textbf{Trucos}
	%----------------------------------------------------------
	
%%%%%%%%%%%%%%%%%%%%%%%%%%%%%%%%%%%%%%%%%%%%%%%%%%%%%%%%%%%
%  SECCIÓN: Pelvis
%  Archivo: secciones/pelvis.tex
%  Sin preámbulo ni \begin{document} — lo maneja main.tex
	%%%%%%%%%%%%%%%%%%%%%%%%%%%%%%%%%%%%%%%%%%%%%%%%%%%%%%%%%%%
	
	\section{Pelvis}
	
	\subsection{Método de Cleaves Modificado (Anca de Rana Bilateral / Lowenstein)}
	
	\subsubsection*{1. Identificación}
	
	\textbf{Tipo:} Complementaria. Forma par con la proyección AP de pelvis/cadera (frente).
	
	\textbf{Denominaciones clínicas frecuentes:} oblicua de cadera/pelvis, perfil de cadera, falso perfil de cadera, Cleaves de cadera.
	
	\textbf{Indicaciones clínicas principales:}
	\begin{itemize}
		\item \textbf{En niños:} displasia del desarrollo de la cadera (DDC), enfermedad de Legg-Calvé-Perthes (necrosis avascular).
		\item \textbf{En adultos:} control de implante o prótesis (seguimiento de osteosíntesis), pinzamiento femoroacetabular, displasia de cadera (no es indispensable visualizar crestas ilíacas en este contexto).
		\item Evaluación comparativa bilateral de coxa vara, coxa valga y epifisiolisis femoral proximal.
	\end{itemize}
	
	\textbf{Ventajas de la toma bilateral:} permite radioprotección al completar el estudio en un solo enfoque y facilita la comparación anatómica simultánea de ambas caderas. \textbf{Desventaja:} el RC no se centra en una región articular precisa sino en el punto medio, lo que puede reducir la definición periférica.
	
	\textbf{Contraindicación absoluta:} sospecha de fractura de cuello femoral, traumatismo agudo de pelvis o luxación de cadera. En estos casos se indica la proyección de Danelius-Miller (rayo horizontal, sin movilización del miembro afectado).
	
	\subsubsection*{2. Técnica Base}
	
	\textbf{Receptor de imagen:} 35 $\times$ 43 cm. Orientación \textbf{transversal} para abarcar ambas caderas simultáneamente.
	
	\textbf{DFRI:} 100 cm. Distancia estándar que equilibra magnificación y dosis.
	
	\textbf{Técnica:}
	\begin{itemize}
		\item \textbf{Paciente joven / neonato:} técnica homogeneizada; prioriza la radioprotección gonadal y permite diferenciar tejidos blandos y hueso con baja densidad calcificada (desarrollo óseo incompleto).
		\item \textbf{Paciente adulto estándar:} técnica intermedia; escala larga de grises para diferenciar hueso, metal (prótesis) y partes blandas, con penetración suficiente para atravesar la densidad de cadera y sacro.
		\item \textbf{Paciente con elevada masa corporal:} técnica básica con Bucky habilitado para garantizar penetración adecuada; si se emplea técnica homogeneizada o intermedia en pacientes de mayor espesor, la radiación dispersa engrisa la imagen (hipodensidad global), por lo que se compensa aumentando mAs para obtener contraste adecuado.
	\end{itemize}
	
	\textbf{Bucky:} Sí. Se utiliza Bucky de mesa con grilla para absorber radiación dispersa generada por el espesor pélvico.
	
	\subsubsection*{3. Posicionamiento}
	
	\paragraph{Paciente}
	Decúbito supino sobre la mesa radiográfica. En pediatría se posiciona al paciente en la punta de la mesa. Plano medio sagital perpendicular al receptor. Manos apoyadas sobre la parte superior del tórax para retirarlas del campo.
	
	\paragraph{Región de interés - Bilateral (anca de rana / Lowenstein)}
	Flexionar caderas y rodillas elevando los muslos hacia el abdomen; llevar los pies hacia arriba buscando que los fémures queden casi verticales (ideal: plantas de los pies juntas como apoyo mutuo). Una vez lograda la flexión, solicitar abducción de ambos muslos aproximadamente 45°, dejando caer las rodillas hacia afuera. Esta abducción sitúa los cuellos femorales paralelos al plano del RI respecto al plano vertical. Verificar que las EIAS se encuentren equidistantes de la mesa para confirmar ausencia de rotación pélvica.
	
	\begin{figure}[h]
		\centering
		\includegraphics[width=0.7\linewidth]{imagenes/pelvis9}
		\caption{}
		\label{fig:pelvis9}
	\end{figure}
	
	
	\paragraph{Región de interés - Unilateral}
	Elevar el hemicuerpo contralateral hasta 45°. El paciente flexiona la cadera y la rodilla del lado afectado y lleva el pie hacia la rodilla contraria todo lo posible, apoyando la planta sobre dicha rodilla. Abducir lateralmente el muslo del lado afectado aproximadamente 45°. El paciente puede inclinarse levemente hacia el lado afectado para facilitar la posición.
	
	\begin{figure}[h]
		\centering
		\includegraphics[width=0.7\linewidth]{imagenes/pelvis10}
		\caption{}
		\label{fig:pelvis10}
	\end{figure}
	
	
	\paragraph{Rayo Central}
	\textbf{Bilateral:} perpendicular al plano del RI. Punto de entrada: plano medio sagital, 2,5 cm por encima de la sínfisis del pubis. Este método (perpendicular) distingue al Cleaves \textit{modificado} del método original, que emplea angulación cefálica.
	
	\textbf{Unilateral:} perpendicular, dirigido hacia el cuello femoral del lado afectado; centrado sobre la EIAS del lado correspondiente.
	
	\paragraph{Receptor de imagen}
	Centrado 2,5 cm por encima de la sínfisis del pubis. Deben incluirse la totalidad de ambas crestas ilíacas y el tercio superior de ambos fémures, con partes blandas adyacentes. En estudios de prótesis o displasia con crestas cortadas, el estudio conserva valor diagnóstico siempre que se visualice la prótesis en toda su extensión.
	
	\paragraph{Respiración}
	Suspendida al momento de la exposición para evitar movimiento de estructuras pélvicas.
	
	\subsubsection*{4. Anatomía Visible}
	
	Cabezas, cuellos y regiones trocantéricas femorales bilaterales proyectadas en forma comparativa. Acetábulos bilaterales en proyección oblicua. Articulaciones coxofemorales en posición de abducción funcional, permitiendo evaluación comparativa del desarrollo articular y de la relación cabeza femoral--acetábulo. Trocánter menor en proyección de perfil sobre el lado medial del fémur. Trocánter mayor superpuesto sobre los cuellos femorales, hacia posterior. Los ejes longitudinales de ambos cuellos femorales deben aparecer paralelos al plano del RI y visibles en toda su extensión entre cabeza y cuerpo femoral.
	
	\subsubsection*{5. Criterios de Evaluación}
	
	\textbf{Posicionamiento correcto:} trocánter menor visible de perfil en el lado medial del fémur; cuello femoral sin acortamiento, visible en toda su extensión (paralelo al RI a 45° de abducción); trocánter mayor superpuesto sobre el cuello hacia posterior; ejes femorales con angulaciones idénticas bilateralmente; prótesis visible en su totalidad cuando corresponda.
	
	\textbf{Ausencia de rotación:} apariencia simétrica de las alas ilíacas y de ambos agujeros obturadores (bilateral); EIAS equidistantes de la mesa.
	
	\textbf{Penetración y contraste:} detalle trabecular visible en cabezas y cuellos femorales; diferenciación clara entre trocánteres mayor y menor; espacios articulares coxofemorales nítidos; diferenciación entre hueso, metal y partes blandas en pacientes con prótesis.
	
	\subsubsection*{6. Puntos Críticos}
	
	\textbf{Exceso de abducción:} el trocánter menor se superpone con el cuerpo del fémur, perdiendo la referencia de perfil; reducir el grado de abducción.
	
	\textbf{Abducción insuficiente:} el trocánter mayor queda por encima del menor y cubre el cuello femoral; el trocánter menor no está de perfil y el cuello aparece acortado; solicitar mayor apertura de muslos hasta los 45° o hasta la tolerancia del paciente.
	
	\textbf{Adaptaciones:} en neonatos utilizar dispositivos de control de abducción y rotación (método Andren-Von Rosén); en niños mayores aceptar variación de 25--45° según tolerancia; en pacientes con osteosíntesis o prótesis, la limitación física o el dolor pueden impedir alcanzar los 45°, lo cual se documenta y se justifica en la solicitud; ante sospecha de fractura de cuello femoral, traumatismo agudo o luxación, está contraindicado movilizar la extremidad: solicitar proyección de Danelius-Miller con rayo horizontal.
	
	\textbf{Trucos:} verificar simetría de abducción palpando ambos trocánteres mayores antes de exponer; asegurar que las plantas de los pies permanezcan juntas como apoyo mutuo para estabilidad de la posición; el nombre ``patas de rana'' describe la similitud visual con la postura del batracio y es útil para comunicar la posición al paciente; la proyección bilateral es preferible cuando la indicación lo permite, ya que permite comparación en una sola exposición y reduce la dosis total.
	
	\subsubsection*{Consideraciones Técnicas por Contexto Clínico}
	
	\textbf{Control de prótesis - paciente pediátrico:} técnica homogeneizada con Bucky habilitado. El metal de la prótesis posee mayor densidad que el hueso y genera radiación dispersa; el Bucky absorbe ese exceso mejorando la nitidez de los bordes metálicos y permitiendo visualizar el pasaje entre hueso y metal. kVp en rango 60--85. Ejemplo de equivalencia para espesor de 16 cm: kVp 75 --- mAs 14 (base); equivalencia ganando penetración para el metal y reduciendo dosis: kVp 86 --- mAs 7.
	
	\textbf{Control de prótesis - paciente adulto:} técnica intermedia.
	
	
	\clearpage
	\onecolumn
	
	\begin{figure}
		\centering
		\includegraphics[width=0.7\linewidth]{imagenes/pelvis11}
		\caption{Método de Cleaves modificado bilateral}
		\label{fig:pelvis11}
	\end{figure}
	
	\begin{figure}
		\centering
		\includegraphics[width=0.7\linewidth]{imagenes/pelvis12}
		\caption{Método de Cleaves modificado bilateral}
		\label{fig:pelvis12}
	\end{figure}
	
	
	\twocolumn
	\clearpage
	
	%----------------------------------------------------------
	%  Para agregar otro enfoque duplicá el bloque
	%  desde \subsection hasta \textbf{Trucos}
	%----------------------------------------------------------
	
	
%%%%%%%%%%%%%%%%%%%%%%%%%%%%%%%%%%%%%%%%%%%%%%%%%%%%%%%%%%%
%  SECCIÓN: Pelvis
%  Archivo: secciones/pelvis.tex
%  Sin preámbulo ni \begin{document} — lo maneja main.tex
	%%%%%%%%%%%%%%%%%%%%%%%%%%%%%%%%%%%%%%%%%%%%%%%%%%%%%%%%%%%
	
	\section{Pelvis}
	
	\subsection{Método de Danelius-Miller - Proyección Axiolateral - (Cross Table / Perfil Quirúrgico)}
	
	\subsubsection*{1. Identificación}
	
	\textbf{Tipo:} Proyección complementaria. Enfoque traumático fundamental; se indica cuando el paciente no puede o no debe realizar la abducción ni ser movilizado.
	
	\textbf{Indicaciones clínicas principales:}
	\begin{itemize}
		\item Fractura de cuello femoral: permite evaluar desplazamiento anterior o posterior, que el frente no evidencia.
		\item Control postoperatorio de cadera y fémur proximal: verifica la correcta colocación del material de osteosíntesis (gamma, DHS, DCS) y evaluación de prótesis de cadera (infección o aflojamiento).
		\item Luxación de cadera y fractura de cabeza femoral.
		\item Fractura de fémur proximal en contexto de traumatismo o politraumatismo.
		\item Pinzamiento femoroacetabular (dolor y restricción de movimiento) y artrosis de cadera.
		\item Necrosis avascular de cabeza femoral y fracturas por estrés del cuello femoral.
		\item Cualquier patología de cadera en la que el paciente no puede ser movilizado (UCI, quirófano, sala de urgencias).
	\end{itemize}
	
	\subsubsection*{2. Técnica Base}
	
	\textbf{Receptor de imagen:} 35 $\times$ 43 cm. Orientación \textbf{longitudinal} respecto al eje mayor del cuello femoral. El RI se coloca vertical, paralelo al cuello femoral (aproximadamente 45° respecto al plano de la mesa), con el borde superior por encima de la cresta ilíaca. El paciente puede sujetar el RI con la mano contralateral; en caso contrario, mantenerlo con sacos de arena.
	
	\textbf{DFRI:} 100 cm. Distancia estándar para técnica sin Bucky; permite compensar la geometría del haz horizontal y mantener una magnificación aceptable.
	
	\textbf{Técnica:} Técnica intermedia. El rayo central atraviesa transversalmente el muslo, región con gran masa muscular y tejidos blandos; se requiere escala de grises amplia para visualizar el trabeculado óseo del fémur proximal y detectar posibles fracturas. Nota: algunos apuntes de clase la consignan como técnica básica ante la ausencia de Bucky y el elevado espesor (mayor dispersión sin rejilla); considerar biotipo del paciente al seleccionar los factores.
	
	\textbf{Bucky:} Sin Bucky. No se utiliza rejilla ni Bucky de mesa ni mural; la ausencia de grilla es consecuencia directa de la geometría del enfoque (RI vertical apoyado en el paciente). No se debe compensar con rejilla.
	
	\subsubsection*{3. Posicionamiento}
	
	\paragraph{Paciente}
	Decúbito supino sobre la mesa radiológica. El plano medio sagital permanece perpendicular a la mesa. El lado afectado se ubica cerca del borde de la mesa para facilitar la colocación del RI. No debe existir rotación pélvica: las EIAS deben ser equidistantes a la superficie de la mesa.
	
	\begin{figure}
		\centering
		\includegraphics[width=0.7\linewidth]{imagenes/pelvis13}
		\caption{}
		\label{fig:pelvis13}
	\end{figure}
	
	\begin{figure}[h]
		\centering
		\includegraphics[width=0.7\linewidth]{imagenes/pelvis14}
		\caption{}
		\label{fig:pelvis14}
	\end{figure}
	
	
	\paragraph{Región de interés}
	Flexionar la rodilla y la cadera del lado \textbf{no afectado} hasta aproximadamente 90°, elevando el muslo a posición vertical; apoyar la pierna sobre un soporte adecuado que no interfiera con el rayo central. El paciente no debe apoyar el pie sobre el tubo de rayos X ni sobre el colimador. La cadera afectada permanece en extensión. De no estar contraindicado por trauma (fractura no consolidada), sujetar el talón del lado afectado y realizar rotación interna de 15--20°. En pacientes delgados o en camas blandas, elevar la pelvis con un almohadón firme para centrar el trocánter mayor en el punto medio del RI. Los repères anatómicos de referencia son: crestas ilíacas, EIAS, pubis y trocánter mayor. Para localizar el trocánter mayor (subcutáneo) se le solicita al paciente que junte la punta de los pies y separe los talones.
	
	\paragraph{Rayo Central}
	Horizontal y perpendicular tanto al RI como al eje longitudinal del cuello femoral. Punto de entrada: \textbf{cara interna del muslo del lado afectado, a nivel del pliegue inguinal}, dirigido hacia el cuello femoral del lado afectado. El haz atraviesa el muslo de medial a lateral.
	
	\paragraph{Receptor de imagen}
	Centrado en el cuello femoral. Borde superior por encima de la cresta ilíaca. La colimación debe abarcar completamente la articulación coxofemoral (acetábulo y cabeza femoral), cuello femoral, ambos trocánteres y el $\frac{1}{2}$ a $\frac{1}{3}$ superior de la diáfisis femoral; debe incluirse además la tuberosidad isquiática (porción del isquion). El RI se mantiene perpendicular al rayo central, sostenido en posición vertical contra el flanco del paciente.
	
	\paragraph{Respiración}
	Suspendida (apnea durante la exposición). La respiración no afecta directamente la región de interés, pero la apnea evita el movimiento del paciente, crítico en politraumatizados y en técnicas con mayor tiempo de exposición (equipos portátiles).
	
	\subsubsection*{4. Anatomía Visible}
	
	Articulación coxofemoral en proyección axiolateral: cabeza femoral en vista lateral verdadera, acetábulo en vista lateral. Cuello femoral desplegado sin superposición del trocánter mayor (criterio fundamental de la proyección). Trocánteres mayor y menor en su relación lateral específica: con rotación interna de 15--20°, el trocánter mayor se proyecta hacia anterior con mínima superposición al cuello, y el trocánter menor queda oculto o levemente superpuesto a la diáfisis femoral hacia posterior. Sin rotación interna (contraindicada por sospecha de fractura), el cuello femoral se verá acortado, aunque aún útil para evaluar desplazamiento lateral. Tuberosidad isquiática proyectada por debajo de la cabeza femoral (referencia para lateralidad: permite identificar si la imagen corresponde a cadera derecha o izquierda; no puede estar cortada). Fémur proximal en vista lateral con posibilidad de evaluar desplazamientos fracturarios. Cualquier dispositivo ortopédico (gamma, DHS, DCS, prótesis) debe verse en su totalidad.
	
	\subsubsection*{5. Criterios de Evaluación}
	
	\textbf{Posicionamiento correcto:} Cuello femoral claramente visible sin superposición del trocánter mayor. Pequeña parte del trocánter menor visible en la superficie posterior del fémur. Con rotación interna adecuada, pequeña porción del trocánter mayor en superficies anterior y posterior. Tuberosidad isquiática proyectada debajo de la cabeza femoral y sin cortar. Todo dispositivo ortopédico incluido en su integridad. RI en posición vertical con orientación correcta (borde superior sobre cresta ilíaca, paralelo al cuello femoral).
	
	\textbf{Ausencia de rotación:} Articulación coxofemoral y acetábulo claramente definidos. Ausencia de sombra de partes blandas del muslo no afectado (glúteos contralaterales) superpuesta sobre la articulación de cadera o el cuello femoral; si el paciente rota el tronco, dichas sombras invaden la zona de interés. Proyección lateral verdadera del fémur proximal.
	
	\textbf{Penetración y contraste:} Detalle trabecular visible en cabeza y cuello femoral. Visualización clara de líneas de fractura si están presentes. Dispositivos ortopédicos con densidad óptica que permita evaluar su integridad completa. Escala de grises adecuada para diferenciar cortical, trabeculado y partes blandas periarticulares.
	
	\subsubsection*{6. Puntos Críticos}
	
	\textbf{Superposición del muslo contralateral:} Principal problema técnico del enfoque. Se produce cuando la flexión de la cadera no afectada es insuficiente o cuando la pierna contralateral no se apoya adecuadamente. Consecuencia: el muslo no afectado bloquea el rayo central o superpone sombra de partes blandas sobre el cuello femoral. Corrección: flexionar la cadera no afectada hasta 90° y asegurar que el soporte no interfiera con el haz; verificar que el paciente no apoye el pie sobre el tubo.
	
	\textbf{Rotación insuficiente del lado afectado:} Sin los 15--20° de rotación interna (cuando no está contraindicada), el trocánter mayor se superpone al cuello femoral ocultando posibles fracturas o desplazamientos. Corrección: sujetar el talón y rotar internamente la extremidad afectada antes de exponer. Si la rotación está contraindicada por fractura, documentar la limitación y evaluar igual el desplazamiento lateral disponible.
	
	\textbf{Adaptaciones:} En fracturas no consolidadas, la manipulación del paciente debe ser realizada exclusivamente por el médico. Si no se puede flexionar la cadera contralateral, considerar el método de Clements-Nakayama (angulación de 15° posterior del RC). En pacientes delgados o en camas blandas, elevar la pelvis con soporte firme para centrar el trocánter mayor en el RI. Con dispositivos ortopédicos, incluir completamente el implante en la imagen. En UCI o quirófano, utilizar equipos portátiles con técnicas adaptadas al entorno.
	
	\textbf{Trucos:} Localizar el trocánter mayor palpándolo (subcutáneo) pidiendo al paciente que junte las puntas de los pies y separe los talones. La tuberosidad isquiática proyectada debajo de la cabeza femoral es la referencia clave para determinar lateralidad (derecha/izquierda) y no debe quedar fuera del campo. Un frente de pelvis solo no es suficiente en traumatismo de cadera: este enfoque axiolateral es imprescindible para evaluar desplazamiento anterior o posterior del cuello femoral. La imagen mal posicionada (RI no vertical o no en decúbito supino) impide determinar correctamente la lateralidad y la relación espacial de los fragmentos.
	
	\clearpage
	\onecolumn
	
	\begin{figure}
		\centering
		\includegraphics[width=0.7\linewidth]{imagenes/pelvis15}
		\caption{}
		\label{fig:pelvis15}
	\end{figure}
	
	\begin{figure}
		\centering
		\includegraphics[width=0.7\linewidth]{imagenes/pelvis16}
		\caption{}
		\label{fig:pelvis16}
	\end{figure}
	
	
	\twocolumn
	\clearpage
	%----------------------------------------------------------
	%  Para agregar otro enfoque duplicá el bloque
	%  desde \subsection hasta \textbf{Trucos}
	%----------------------------------------------------------
	
	
%%%%%%%%%%%%%%%%%%%%%%%%%%%%%%%%%%%%%%%%%%%%%%%%%%%%%%%%%%%
%  SECCIÓN: Pelvis
%  Archivo: secciones/pelvis.tex
%  Sin preámbulo ni \begin{document} — lo maneja main.tex
	%%%%%%%%%%%%%%%%%%%%%%%%%%%%%%%%%%%%%%%%%%%%%%%%%%%%%%%%%%%
	
	\section{}
	
	\subsection{Axiolateral Modificada - Método de Clements-Nakayama}
	
	\subsubsection*{1. Identificación}
	
	\textbf{Tipo:} Complementario. Alternativa al método de Danelius-Miller (cross table) cuando no es posible realizarlo.
	
	\textbf{Indicaciones clínicas principales:}
	\begin{itemize}
		\item Fracturas bilaterales de cadera: método de elección cuando el paciente no puede elevar ninguna pierna
		\item Artroplastia bilateral de cadera: seguimiento post-quirúrgico sin stress adicional sobre los implantes
		\item Politraumatismo grave con inmovilización total: mínima manipulación del paciente crítico (UCI)
		\item Limitación severa de movilidad bilateral: artrosis avanzada, rigidez articular o cualquier condición que impida flexionar la cadera contralateral
	\end{itemize}
	
	\subsubsection*{2. Técnica Base}
	
	\textbf{Receptor de imagen:} 24 x 30 cm. Orientación \textbf{longitudinal} (vertical), apoyado en la cara lateral del muslo sobre la cresta ilíaca, inclinado 15° hacia posterior para que el RC incida perpendicular al RI.
	
	\textbf{DFRI:} 100 cm. Distancia estándar para proyecciones de cadera.
	
	\textbf{Técnica:} Intermedia. Se requiere visualizar el trabeculado óseo del fémur con un kVp intermedio que compense la larga trayectoria del rayo a través de la masa muscular y los tejidos blandos de la región. En presencia de prótesis, ajustar para evitar que el implante genere una sombra hiperdensa que tape estructuras adyacentes.
	
	\textbf{Bucky:} No se utiliza Bucky de mesa. Se emplea una rejilla antidifusora portátil con líneas paralelas al suelo, posicionada paralela al eje del cuello femoral.
	
	\subsubsection*{3. Posicionamiento}
	
	\paragraph{Paciente}
	Decúbito supino estricto sobre la mesa radiográfica, con el lado afectado cerca del borde de la mesa. No se realiza ningún movimiento de las extremidades inferiores: el miembro se mantiene en la posición neutral con la que llega el paciente.	
	
	\paragraph{Región de interés}
	El miembro inferior afectado permanece en posición neutral, sin rotación interna ni externa forzada. No se manipula la extremidad. La pelvis debe estar sin rotación.
	
	% TODO: \usepackage{graphicx} required
	\begin{figure}[h]
		\centering
		\includegraphics[width=0.7\linewidth]{imagenes/pelvis17}
		\caption{}
		\label{fig:pelvis17}
	\end{figure}
	
	
	\paragraph{Rayo Central}
	Angulado 15° en sentido posterior, alineado perpendicular al cuello femoral y al RI. Punto de entrada: \textbf{cara interna del muslo}, dirigido hacia la articulación coxofemoral. La angulación posterior de 15° es fundamental para esquivar el muslo contralateral y obtener un verdadero perfil lateral.
	
		% TODO: \usepackage{graphicx} required
	\begin{figure}[h]
		\centering
		\includegraphics[width=0.7\linewidth]{imagenes/pelvis18}
		\caption{}
		\label{fig:pelvis18}
	\end{figure}
	
	\paragraph{Receptor de imagen}
	Apoyado en la cara lateral del muslo, con su borde inferior por debajo del paciente y el borde posterior angulado 15° hacia posterior. Centrado sobre la articulación coxofemoral. Debe incluir: acetábulo, cabeza femoral, cuello femoral, trocánteres mayor y menor en perfil lateral. La inclusión de la \textbf{tuberosidad isquiática} es indicador directo de que el RI quedó con la inclinación correcta de 15°: si la tuberosidad aparece cortada, significa que el RI quedó vertical a 90° y la imagen se proyectó fuera del receptor.
	
	% TODO: \usepackage{graphicx} required
	\begin{figure}[h]
		\centering
		\includegraphics[width=0.7\linewidth]{imagenes/pelvis19}
		\caption{}
		\label{fig:pelvis19}
	\end{figure}
	
	
	\paragraph{Respiración}
	Suspendida durante la exposición.
	
	\subsubsection*{4. Anatomía Visible}
	
	Se visualiza la articulación coxofemoral en perfil lateral verdadero: cabeza femoral dentro del acetábulo, cuello femoral (típicamente acortado en reposo o por fractura), trocánter mayor superpuesto al cuello femoral, y trocánter menor de perfil hacia posterior. Se debe incluir el fémur proximal completo. En pacientes con dispositivos ortopédicos, el implante debe verse en su integridad. La tuberosidad isquiática debe estar presente en la imagen como referencia de correcta angulación del RI. A diferencia del método cross table, la tuberosidad isquiática puede aparecer con mayor distorsión debido a la angulación posterior del RC.
	
	\subsubsection*{5. Criterios de Evaluación}
	
	\textbf{Posicionamiento correcto:} Articulación coxofemoral claramente demostrada en perfil lateral. Cabeza, cuello y trocánteres femorales en perfil lateral. Trocánter mayor superpuesto al cuello femoral. Trocánter menor visible de perfil hacia posterior. Tuberosidad isquiática incluida en la imagen. Dispositivos ortopédicos visibles en su totalidad.
	
	\textbf{Ausencia de rotación:} Ausencia de sombra densa sobre el cuello femoral proveniente del muslo contralateral. Si se observa superposición de la pierna sana tapando el cuello, indica que el RC no se angulò los 15° suficientes para esquivar el muslo contralateral.
	
	\textbf{Penetración y contraste:} Visualización del trabeculado óseo del fémur proximal. Contraste adecuado para evaluar el espacio articular y, cuando corresponda, la integridad de implantes ortopédicos. Técnica intermedia que equilibre resolución ósea con compensación de tejidos blandos.
	
	\subsubsection*{6. Puntos Críticos}
	
	\textbf{Angulación insuficiente del RC:} Si el RC no alcanza los 15° posteriores exactos, el muslo contralateral genera una sombra densa que se superpone al cuello femoral, comprometiendo el diagnóstico. Corrección: verificar la angulación del tubo antes de exponer.
	
	\textbf{Tuberosidad isquiática ausente o cortada:} Indica que el RI quedó vertical (90°) en lugar de inclinado 15°. El RC al entrar angulado proyecta la imagen fuera del receptor. Una placa de Clements-Nakayama sin tuberosidad isquiática se considera no diagnóstica en traumatismos de alta energía, ya que no permite descartar luxación posterior con seguridad. Corrección: inclinar el borde posterior del RI 15° hacia posterior antes de exponer.
	
	\textbf{Adaptaciones:} Método de elección cuando el Danelius-Miller (cross table) no es posible: fracturas bilaterales, artroplastia bilateral, pacientes en UCI con inmovilización total, o cualquier limitación que impida flexionar la cadera contralateral. Requiere rejilla portátil con capacidad de angulación; no utiliza Bucky de mesa.
	
	\textbf{Trucos:} La ventaja principal sobre el Danelius-Miller es que no requiere flexionar la pierna sana, lo que evita superposición del muslo contralateral y reduce la manipulación del paciente. Los 15° posteriores del RC son el parámetro crítico de todo el enfoque: determinan tanto la calidad diagnóstica de la imagen como la correcta proyección de la tuberosidad isquiática dentro del RI. Verificar siempre que las líneas de la rejilla queden paralelas al suelo y que el RI esté paralelo al eje del cuello femoral para evitar distorsión geométrica.
	
	
	\onecolumn
	\clearpage
	
	% TODO: \usepackage{graphicx} required
	\begin{figure}
		\centering
		\includegraphics[width=0.7\linewidth]{imagenes/pelvis20}
		\caption{}
		\label{fig:pelvis20}
	\end{figure}
	
	
	\twocolumn
	\clearpage
	%----------------------------------------------------------
	%  Para agregar otro enfoque duplicá el bloque
	%  desde \subsection hasta \textbf{Trucos}
	%----------------------------------------------------------
	
%%%%%%%%%%%%%%%%%%%%%%%%%%%%%%%%%%%%%%%%%%%%%%%%%%%%%%%%%%%
%  SECCIÓN: Pelvis
%  Archivo: secciones/pelvis.tex
%  Sin preámbulo ni \begin{document} — lo maneja main.tex
	%%%%%%%%%%%%%%%%%%%%%%%%%%%%%%%%%%%%%%%%%%%%%%%%%%%%%%%%%%%
	
	\section{Pelvis}
	
	\subsection{Lequesne - Falso Perfil Parado}
	
	\subsubsection*{1. Identificación}
	
	\textbf{Tipo:} Complementario. Forma par con AP de pelvis (frente) en displasia de cadera y coxartrosis. En pinzamiento femoroacetabular el par es AP de pelvis y método de Cleaves, siendo el Lequesne un complementario opcional junto al Dunn 45° o 90°.
	
	\textbf{Indicaciones clínicas principales:}
	\begin{itemize}
		\item Displasia de cadera del adulto: evaluación del grado de cobertura anterior de la cabeza femoral en carga
		\item Coxartrosis / artrosis de cadera: valoración del espacio articular coxofemoral con soporte de peso
		\item Pinzamiento femoroacetabular (FAI): estudio complementario para ver la transición cuello-cabeza
		\item Sospecha de rotura labral o condral (rodete acetabular): visualización del reborde acetabular anterior
		\item Evaluación pre y post-quirúrgica de osteotomías periacetabulares
	\end{itemize}
	
	\subsubsection*{2. Técnica Base}
	
	\textbf{Receptor de imagen:} 35 × 43 cm. Orientación \textbf{transversal (horizontal / apaisado)} para incluir ambas articulaciones coxofemorales y el tercio proximal de fémur.
	
	\textbf{DFRI:} 100 cm. Distancia estándar de 1 metro; la magnificación resultante es intencional y permite apreciar con mayor detalle el grado de cobertura acetabular sobre la cabeza femoral.
	
	\textbf{Técnica:} Intermedia con alto kVp. El alto kVp compensa el mayor espesor que genera la rotación del tronco (65°), que superpone mayor cantidad de tejido óseo y blando en el trayecto del rayo.
	
	\textbf{Bucky:} Bucky mural o de pie; utilizar rejilla antidifusora obligatoria dado el espesor de la región.
	
	\subsubsection*{3. Posicionamiento}
	
	\paragraph{Paciente}
	De pie, descalzo, con los brazos cruzados a la altura del pecho. La bipedestación es condición indispensable: al soportar su propio peso, la evaluación de la cobertura acetabular es funcional. Si existe diferencia de longitud en miembros inferiores, colocar apoyo bajo el pie más corto para igualar la distribución del peso. El paciente puede agarrarse al soporte lateral del bucky para estabilidad, siempre que no altere la posición.
	
	% TODO: \usepackage{graphicx} required
	\begin{figure}[h]
		\centering
		\includegraphics[width=0.7\linewidth]{imagenes/pelvis22}
		\caption{}
		\label{fig:pelvis22}
	\end{figure}
	
	
	\paragraph{Región de interés}
	Partiendo de un perfil estricto (90°) de la cadera afectada con el PMS paralelo al receptor, se rota el cuerpo 25° hacia atrás, formando un ángulo de \textbf{65°} respecto al plano del receptor. El pie de la cadera afectada queda \textbf{paralelo al receptor} (esto ubica el fémur en perfil estricto); el pie contralateral queda \textbf{perpendicular al receptor} (mirando hacia el tubo). La cadera afectada es la que queda apoyada contra el receptor.
	
	\paragraph{Rayo Central}
	Perpendicular al plano del receptor. Punto de entrada: \textbf{2 cm por debajo del punto medio del pliegue inguinal} de la cadera estudiada (cadera apoyada). Algunas referencias bibliográficas ubican el centro aproximadamente 5 cm por encima del trocánter mayor; ambas referencias apuntan a la articulación coxofemoral apoyada.
	
	\paragraph{Receptor de imagen}
	Centrado en la articulación coxofemoral de la cadera afectada. La inclusión mínima debe abarcar: hacia superior, la parte inferior del ala ilíaca; hacia distal, el tercio proximal del fémur con el trocánter menor visible; sínfisis pubiana y cadera contralateral; partes blandas hacia lateral externo. La cadera contralateral es imprescindible para medir la distancia intercoxofemoral y validar la rotación.
	
	\paragraph{Respiración}
	Suspendida en espiración. Reduce el movimiento de la cintura pélvica durante la exposición y asegura la reproducibilidad de la posición.
	
	\subsubsection*{4. Anatomía Visible}
	
	Se visualiza la articulación coxofemoral de la cadera apoyada en perfil, con el reborde anterior del acetábulo proyectado por delante de la cabeza femoral. La articulación coxofemoral contralateral (alejada) debe aparecer en la imagen con el agujero obturador abierto por la oblicuidad, necesario para medir la distancia intercoxofemoral. En la cadera apoyada el agujero obturador se presenta estrecho o casi cerrado por estar de perfil. Son visibles la pared posterior del acetábulo, el espacio articular medial e inferior (zona de contragolpe), el sacro, el cóccix y la sínfisis pubiana. Los fémures proximales aparecen superpuestos entre sí, al igual que las sombras acetabulares y los márgenes posteriores del isquion e ilion.
	
	\subsubsection*{5. Criterios de Evaluación}
	
	\textbf{Posicionamiento correcto:} La distancia entre las dos articulaciones coxofemorales debe ser equivalente a \textbf{una cabeza y media femoral}. Si la distancia supera 1,5 cabezas femorales, hubo exceso de rotación (el paciente giró demasiado hacia el frente); si es menor a 1,5 cabezas femorales, faltó rotación (el paciente debe rotar más hacia atrás). La cadera apoyada debe mostrar el espacio articular despejado y el reborde acetabular nítido. El reborde anterior del acetábulo debe verse proyectado por delante de la cabeza femoral. El trocánter menor debe ser visible.
	
	\textbf{Ausencia de rotación:} El pie de la cadera apoyada paralelo al receptor indica correcta orientación del fémur en perfil. El agujero obturador de la cadera apoyada cerrado y el de la cadera contralateral abierto confirman la oblicuidad de 65°.
	
	\textbf{Penetración y contraste:} Adecuada visualización de los espacios articulares y contornos acetabulares con nitidez ósea. El alto kVp debe compensar el espesor sin sacrificar la definición del reborde acetabular ni del espacio articular medial-inferior.
	
	\subsubsection*{6. Puntos Críticos}
	
	\textbf{Rotación incorrecta (no exactamente 65°):} Una rotación insuficiente o excesiva modifica la distancia intercoxofemoral más o menos allá de 1,5 cabezas femorales, impidiendo la correcta evaluación de la cobertura anterior. Controlar siempre la posición del pie apoyado (paralelo al RI) como referencia directa del ángulo femoral.
	
	\textbf{Pie no paralelo al receptor:} Si el pie de la cadera apoyada no queda paralelo al receptor, el fémur no queda en perfil estricto y el reborde acetabular anterior se distorsiona, invalidando la medición de cobertura. Es el error más frecuente y el más fácil de prevenir revisando la posición del pie antes de exponer.
	
	\textbf{Adaptaciones:} Contraindicado en trauma agudo (sospecha de fractura o luxación de cadera) y en enfermedad destructiva activa de cadera, ya que la bipedestación y la rotación forzada pueden agravar la lesión. En pacientes con limitaciones físicas proporcionar apoyo lateral y verificar distribución equitativa del peso.
	
	\textbf{Trucos:} La proyección es funcional y comparativa: se realiza con carga y frecuentemente de ambas caderas. El marcador lateral debe identificar claramente el lado estudiado, ya que la imagen se orienta y marca como un frente (incluye ambas coxofemorales, el marcador queda hacia el lateral del muslo). Tomar la referencia de 1,5 cabezas femorales según la bibliografía principal del curso; la variante de 1 cabeza existe en otros textos. El método de Cleaves no reemplaza al Lequesne en displasia porque se realiza en decúbito y no evalúa la cadera en carga.
	
	\onecolumn
	\clearpage
	
	% TODO: \usepackage{graphicx} required
	\begin{figure}
		\centering
		\includegraphics[width=0.7\linewidth]{imagenes/pelvis21}
		\caption{}
		\label{fig:pelvis21}
	\end{figure}
	
	\twocolumn
	\clearpage
	%----------------------------------------------------------
	%  Para agregar otro enfoque duplicá el bloque
	%  desde \subsection hasta \textbf{Trucos}
	%----------------------------------------------------------
	
%%%%%%%%%%%%%%%%%%%%%%%%%%%%%%%%%%%%%%%%%%%%%%%%%%%%%%%%%%%
%  SECCIÓN: Pelvis
%  Archivo: secciones/pelvis.tex
%  Sin preámbulo ni \begin{document} — lo maneja main.tex
	%%%%%%%%%%%%%%%%%%%%%%%%%%%%%%%%%%%%%%%%%%%%%%%%%%%%%%%%%%%
	
	\section{Pelvis}
	
	\subsection{Método de Teufel (Acetábulo) - OAD / OAI}
	
	\subsubsection*{1. Identificación}
	
	\textbf{Tipo:} Complementaria. Forma par con AP de pelvis con rotación interna estricta y método de Cleaves unilateral.
	
	\textbf{Indicaciones clínicas principales:}
	\begin{itemize}
		\item Fractura acetabular, especialmente del borde supero-posterior del acetábulo
		\item Evaluación de la fóvea capitis del fémur (ligamento redondo)
		\item Luxaciones de cadera y displasia acetabular
	\end{itemize}
	
	\subsubsection*{2. Técnica Base}
	
	\textbf{Receptor de imagen:} 24 $\times$ 30 cm. Orientación \textbf{longitudinal}.
	
	\textbf{DFRI:} 100 cm. Distancia estándar para proyecciones de cadera.
	
	\textbf{Técnica:} Intermedia. Se requiere detalle suficiente para diferenciar la fóvea capitis (pequeña depresión en la cabeza femoral) de la densidad ósea del acetábulo circundante.
	
	\textbf{Bucky:} Bucky de mesa con rejilla.
	
	\subsubsection*{3. Posicionamiento}
	
	\paragraph{Paciente}
	Decúbito prono. Se apoya el lado afectado sobre la mesa. Se coloca almohada bajo la cabeza para mayor comodidad. El cuerpo se eleva sobre el antebrazo del lado afectado y el paciente flexiona la rodilla del lado elevado para mantener la posición estable.
	
	% TODO: \usepackage{graphicx} required
	\begin{figure}[h]
		\centering
		\includegraphics[width=0.7\linewidth]{imagenes/pelvis23}
		\caption{}
		\label{fig:pelvis23}
	\end{figure}
	
	
	\paragraph{Región de interés}
	Se eleva el tórax y la pelvis del lado sano entre 35° y 40° (referencia: 38°) respecto al plano de la mesa, creando una oblicuidad anterior. El lado afectado permanece apoyado. La cadera afectada se centra con la línea media de la rejilla. La proyección es unilateral: no es necesario realizar ambas oblicuas, únicamente la del lado afectado.
	
	\paragraph{Rayo Central}
	Angulado \textbf{cefálico 12°} respecto al plano del receptor. La angulación compensa los 38° de rotación corporal y permite que el RC incida de forma tangencial sobre la concavidad de la fóvea capitis. Punto de entrada: \textbf{cara posterior, a nivel del trocánter mayor del lado apoyado, aproximadamente 5 cm lateral al plano medio sagital hacia el lado en estudio}; el rayo penetra a un nivel inferior al cóccix y sale por el centro del acetábulo.
	
	\paragraph{Receptor de imagen}
	Centrado en la articulación coxofemoral del lado afectado. Debe incluir: parte inferior del ala ilíaca, reborde superior del acetábulo, cabeza femoral completa y partes blandas laterales de la cadera.
	
	\paragraph{Respiración}
	Suspendida durante la exposición para evitar movimiento y asegurar nitidez en las estructuras óseas finas como la fóvea capitis.
	
	\subsubsection*{4. Anatomía Visible}
	
	Acetábulo en proyección axial oblicua. Cabeza femoral de perfil. Pared y borde supero-posterior del acetábulo libres de superposición con la cabeza femoral. Fóvea capitis visible como área cóncava en la cabeza femoral, resultado de la correcta oblicuidad (35--40°) y la angulación cefálica del RC, que permite visualizarla sin que quede cubierta por el pubis o el isquion. Articulación coxofemoral magnificada por ser el lado más cercano al receptor. Trocánter mayor de perfil hacia lateral. Agujero obturador abierto y amplio. Relación entre cabeza femoral y acetábulo claramente definida.
	
	\subsubsection*{5. Criterios de Evaluación}
	
	\textbf{Posicionamiento correcto:} Articulación coxofemoral afectada cercana al centro de la radiografía y magnificada. Borde supero-posterior del acetábulo libre de superposición con la cabeza femoral. Trocánter mayor de perfil. Agujero obturador abierto y amplio. Fóvea capitis visible como concavidad sin superposición del pubis ni isquion.
	
	\textbf{Ausencia de rotación:} Cabeza femoral en verdadero perfil. Agujero obturador proyectado abierto (ni colapsado ni en perfil cerrado). Anatomía acetabular claramente delimitada sin distorsión por exceso ni defecto de oblicuidad.
	
	\textbf{Penetración y contraste:} Detalle óseo óptimo del acetábulo con contraste suficiente para diferenciar la fóvea capitis de la densidad general de la cabeza femoral. Ausencia de sobreexposición o subexposición que comprometa la lectura de estructuras finas.
	
	\subsubsection*{6. Puntos Críticos}
	
	\textbf{Falta de angulación o rotación insuficiente (menos de 35°):} La fóvea no se verá cóncava porque el RC no incide de forma tangencial sobre la depresión. Se corrige verificando los 12° de angulación cefálica y los 35--38° de elevación del lado sano.
	
	\textbf{Exceso de rotación:} El agujero obturador se ve de perfil (cerrado), indicando que el cuerpo giró demasiado. Se corrige bajando el lado elevado hasta recuperar los 38°.
	
	\textbf{Falta de rotación (cuerpo muy de frente):} La cabeza femoral queda tapada por el isquion. Se corrige subiendo el lado elevado para aumentar la oblicuidad.
	
	\textbf{Adaptaciones:} En pacientes con trauma, evitar movimientos forzados; usar apoyos adicionales para mantener los 38° sin comprometer la integridad del paciente. En pacientes con limitaciones de movilidad, incrementar el soporte bajo el tórax y la pelvis del lado sano.
	
	\textbf{Trucos:} Palpar el trocánter mayor antes de centrar para asegurar el punto de entrada correcto del RC. Verificar que el ángulo de 38° se mantenga constante durante toda la exposición usando un goniómetro o cuña radiolúcida. Confirmar visualmente que el agujero obturador se proyecta abierto antes de exponer, como indicador rápido de oblicuidad adecuada.
	
	\onecolumn
	\clearpage
	
	% TODO: \usepackage{graphicx} required
	\begin{figure}
		\centering
		\includegraphics[width=0.7\linewidth]{imagenes/pelvis24}
		\caption{}
		\label{fig:pelvis24}
	\end{figure}
	
	
	\twocolumn
	\clearpage
	
	
	
	%----------------------------------------------------------
	%  Para agregar otro enfoque duplicá el bloque
	%  desde \subsection hasta \textbf{Trucos}
	%----------------------------------------------------------
	
%%%%%%%%%%%%%%%%%%%%%%%%%%%%%%%%%%%%%%%%%%%%%%%%%%%%%%%%%%%
%  SECCIÓN: Pelvis
%  Archivo: secciones/pelvis.tex
%  Sin preámbulo ni \begin{document} — lo maneja main.tex
	%%%%%%%%%%%%%%%%%%%%%%%%%%%%%%%%%%%%%%%%%%%%%%%%%%%%%%%%%%%
	
	\section{Pelvis}
	
	\subsection{Outlet / Proyección AP Axial de Salida - Método de Taylor}
	
	\subsubsection*{1. Identificación}
	
	\textbf{Tipo:} Complementaria. Proyección especializada del anillo pélvico anterior.
	
	\textbf{Indicaciones clínicas principales:}
	\begin{itemize}
		\item Traumatismo pélvico de alta energía con sospecha de fracturas en ramas pubianas e isquiopubianas
		\item Fractura de libro abierto y desplazamientos supero-inferiores de la sínfisis del pubis
		\item Diástasis de sínfisis púbica y lesiones del arco anterior pélvico
		\item Fracturas por estrés en huesos púbicos y lesiones deportivas del arco anterior
		\item Evaluación del anillo pélvico anterior en politraumatismo
	\end{itemize}
	
	\subsubsection*{2. Técnica Base}
	
	\textbf{Receptor de imagen:} 35$\times$43 cm. Orientación \textbf{transversal} para incluir el anillo pélvico anterior completo. Se acepta 24$\times$30 cm cuando el campo de interés es restringido.
	
	\textbf{DFRI:} 100 cm. Distancia estándar para pelvis; el rayo angulado requiere compensación del centro del RI según la proyección del RC sobre la mesa.
	
	\textbf{Técnica:} Intermedia. El rayo debe atravesar volumen considerable de partes blandas (glúteos y pelvis), pero las ramas pubianas no son tan densas como el sacro; la técnica intermedia equilibra detalle trabecular con penetración adecuada del tejido blando. Se ajusta según biotipo del paciente.
	
	\textbf{Bucky:} Bucky de mesa con grilla. El volumen de partes blandas justifica el uso de grilla para reducir radiación dispersa.
	
	\subsubsection*{3. Posicionamiento}
	
	\paragraph{Paciente}
	Decúbito supino con el plano medio sagital perpendicular al RI y centrado sobre la línea media de la grilla. Brazos en posición cómoda, alejados del área de interés.
	
	\begin{figure}[h]
		\centering
		\includegraphics[width=0.7\linewidth]{imagenes/pelvis28}
		\caption{}
		\label{fig:pelvis28}
	\end{figure}
	
	
	\paragraph{Región de interés}
	EIAS equidistantes a la mesa para asegurar ausencia de rotación pélvica. Miembros inferiores en rotación interna (puntas de pies juntas, talones separados) salvo sospecha de fractura de cadera. Rodillas ligeramente flexionadas con soporte si el paciente lo requiere.
	
	\paragraph{Rayo Central}
	Angulado en dirección \textbf{cefálica}: \textbf{20--35°} en hombres y \textbf{30--45°} en mujeres. Punto de entrada: \textbf{5 cm distal al borde superior de la sínfisis del pubis}. La diferencia de angulación por sexo responde a la distinta versión pélvica: los hombres tienden a la retroversión y las mujeres a la anteversión. El centro del RI se ajusta para coincidir con el punto de proyección del RC sobre la bandeja del Bucky.
	
	\paragraph{Receptor de imagen}
	Centrado de forma que el punto medio del RI coincida con la proyección del rayo central sobre la bandeja. La región a incluir comprende las ramas pubianas superior e inferior, sínfisis del pubis, agujeros obturadores y articulaciones coxofemorales con fémur proximal y trocánteres. Existe variante que incluye pelvis completa con crestas ilíacas.
	
	\paragraph{Respiración}
	Suspendida en apnea en el momento de la exposición para evitar movimiento involuntario de estructuras abdominopélvicas.
	
	\subsubsection*{4. Anatomía Visible}
	
	La proyección despliega el anillo pélvico anterior en visión axial inferior. Se visualizan los huesos púbicos e isquiáticos magnificados (efecto propio de la angulación), las ramas pubianas superiores e inferiores sin acortamiento proyeccional, y las ramas isquiáticas claramente delineadas. La sínfisis del pubis se observa desde abajo, magnificada y con cierta distorsión esperable. Los agujeros obturadores aparecen abiertos y simétricos. Se incluyen las articulaciones coxofemorales con el fémur proximal y, según el RI utilizado, las crestas ilíacas. El sacro y el cóccix quedan superpuestos por los huesos pubianos como consecuencia de la angulación cefálica.
	
	\subsubsection*{5. Criterios de Evaluación}
	
	\textbf{Posicionamiento correcto:} Ramas pubianas superpuestas entre S1 y S2. Huesos pubianos e isquiáticos magnificados y centrados en la radiografía. Sínfisis del pubis visible desde abajo con magnificación esperable. Articulaciones coxofemorales incluidas. Cuerpo y rama del isquion con mínima superposición.
	
	\textbf{Ausencia de rotación:} Agujeros obturadores simétricos en tamaño y forma. Ilion bilateral simétrico. EIAS equidistantes al plano de la mesa. Sínfisis púbica centrada en la línea media.
	
	\textbf{Penetración y contraste:} Detalle trabecular visible en ramas púbicas e isquiáticas. Diferenciación clara entre estructuras óseas superpuestas. Visualización de partes blandas adyacentes con escala de grises adecuada para hueso de densidad intermedia.
	
	\subsubsection*{6. Puntos Críticos}
	
	\textbf{Falta de angulación:} Las ramas pubianas se proyectan por debajo del sacro en lugar de superponerse a S1--S2. Corregir aumentando los grados cefálicos dentro del rango para el sexo del paciente.
	
	\textbf{Exceso de angulación:} Las ramas pubianas se proyectan por encima del promontorio sacro. Corregir reduciendo la angulación cefálica.
	
	\textbf{Adaptaciones:} En trauma de alta energía no movilizar los miembros inferiores a rotación interna si se sospecha fractura de cadera asociada; dejar los MMII en posición neutra. En pacientes con dolor intenso apoyar las rodillas con soporte para reducir tensión lumbar y facilitar la cooperación durante la apnea.
	
	\textbf{Trucos:} El lado que aparece con la estructura más grande en imagen es el lado elevado (mayor DOR), ya que la mayor distancia objeto--receptor produce mayor magnificación. El lado apoyado mostrará la estructura más pequeña o cerrada. Para corregir rotación, verificar que las EIAS estén equidistantes a la mesa. Recordar que la sínfisis del pubis, por ser la estructura más anterior, siempre aparecerá magnificada y con cierto desplazamiento aparente; esto es esperable y no debe confundirse con patología.
	
		\onecolumn
	\clearpage
	
	\begin{figure}
		\centering
		\includegraphics[width=0.7\linewidth]{imagenes/pelvis26}
		\caption{En esta imagen no se observa, pero la Rx final debe de contener toda la pelvis completa visible}
		\label{fig:pelvis26}
	\end{figure}
	
	
	\twocolumn
	\clearpage
	
	%----------------------------------------------------------
	%  Para agregar otro enfoque duplicá el bloque
	%  desde \subsection hasta \textbf{Trucos}
	%----------------------------------------------------------
	
	
	%%%%%%%%%%%%%%%%%%%%%%%%%%%%%%%%%%%%%%%%%%%%%%%%%%%%%%%%%%%
	%  SECCIÓN: Pelvis
	%  Archivo: secciones/pelvis_inlet.tex
	%  Sin preámbulo ni \begin{document} — lo maneja main.tex
		%%%%%%%%%%%%%%%%%%%%%%%%%%%%%%%%%%%%%%%%%%%%%%%%%%%%%%%%%%%
		
		\section{Pelvis}
		
		\subsection{Inlet / Proyección Supero-Inferior Axial de Entrada - Método de Bridgman}
		
		\subsubsection*{1. Identificación}
		
		\textbf{Tipo:} Proyección complementaria (enfoque especializado). Forma par con la proyección Outlet (Supero-Inferior Axial de Salida).
		
		\textbf{Indicaciones clínicas principales:}
		\begin{itemize}
			\item Traumatismo pélvico de alta energía con sospecha de fractura del anillo pélvico posterior
			\item Diastasis (separación) de la sínfisis púbica
			\item Fractura con desplazamiento antero-posterior: evaluación de rotación interna/externa del hueso ilíaco (fractura ``libro abierto'')
			\item Fractura del sacro
			\item Luxaciones sacroilíacas
			\item Control posoperatorio para verificar posición de tornillos
			\item Valoración de la integración del anillo pélvico incluyendo la línea innominada
			\item Análisis de la entrada pélvica en obstetricia
			\item Evaluación de deformidades pélvicas congénitas
		\end{itemize}
		
		\subsubsection*{2. Técnica Base}
		
		\textbf{Receptor de imagen:} 35$\times$43 cm. Orientación \textbf{transversal}. Se acepta también 24$\times$30 cm según disponibilidad.
		
		\textbf{DFRI:} 100 cm (1 metro). Distancia estándar para pelvis, permite una cobertura adecuada del anillo pélvico con la angulación caudal aplicada.
		
		\textbf{Técnica:} Intermedia. Requiere suficiente kVp para atravesar toda la masa del sacro y la pelvis ósea, y proyectar correctamente el anillo pélvico; el alto kV permite además diferenciar estructuras superpuestas.
		
		\textbf{Bucky:} Bucky de mesa con grilla.
		
		\subsubsection*{3. Posicionamiento}
		
		\paragraph{Paciente}
		Decúbito supino sobre la mesa radiográfica, con almohada bajo la cabeza. La línea media sagital debe quedar perpendicular al receptor de imagen y centrada en el mismo, equidistante a ambos bordes laterales de la mesa.
		
		\begin{figure}[h]
			\centering
			\includegraphics[width=0.7\linewidth]{imagenes/pelvis31}
			\caption{}
			\label{fig:pelvis31}
		\end{figure}
		
		
		\paragraph{Región de interés}
		Las EIAS deben quedar equidistantes de la mesa, garantizando ausencia de rotación pélvica. Se flexionan ligeramente las rodillas colocando un soporte bajo las mismas para aliviar la tensión lumbar y pélvica. Los miembros inferiores se posicionan con rotación interna leve. Los brazos se ubican en posición cómoda, alejados del área de exposición. El plano medio sagital debe coincidir con la línea media de la grilla.
		
		\paragraph{Rayo Central}
		Angulado \textbf{40° en sentido caudal} respecto al plano vertical del receptor, tanto en hombres como en mujeres. Punto de entrada: \textbf{a nivel de las EIAS (espinas ilíacas anterosuperiores)}, por encima de la sínfisis púbica. La angulación caudal proyecta el anillo pélvico hacia inferior, acortando el sacro y ocultando el cóccix.
		
		\paragraph{Receptor de imagen}
		Centrado a la altura de los trocánteres mayores, coincidiendo con la sínfisis púbica (el RI se desplaza porque el RC entra angulado). Debe incluir superiormente el promontorio del sacro y la sínfisis con las ramas superiores; lateralmente, más allá de los acetábulos y cabezas femorales (hasta los trocánteres mayores). No es necesario incluir las crestas ilíacas.
		
		\paragraph{Respiración}
		Apnea durante la exposición (respiración suspendida). Evita el movimiento de estructuras pélvicas y abdominales durante la adquisición.
		
		\subsubsection*{4. Anatomía Visible}
		
		Se visualiza el anillo pélvico completo (entrada de la pelvis en su totalidad) desde una perspectiva superior, lo que permite evaluar el estrecho superior de la pelvis. Las ramas superior e inferior de los huesos púbicos quedan superpuestas en la porción medial. Los dos tercios laterales de los huesos pubianos e isquiáticos se presentan en casi superposición. Las espinas isquiáticas aparecen simétricas respecto a la línea media, confirmando centrado correcto del RC en el plano medio sagital. Las articulaciones coxofemorales son visibles. Los huesos pélvicos anteriores se proyectan claramente. Las articulaciones sacroilíacas y la sínfisis púbica son evaluables. El sacro se acorta por la angulación caudal y el cóccix queda superpuesto/oculto. Es una vista desde arriba: se evalúan los diámetros pélvicos y la entrada pélvica verdadera.
		
		\subsubsection*{5. Criterios de Evaluación}
		
		\textbf{Posicionamiento correcto:} Las ramas superior e inferior de los huesos pubianos deben estar superpuestas en la porción medial. Los dos tercios laterales de los huesos pubianos e isquiáticos en casi superposición. El anillo pélvico debe mostrar morfología circular o con forma de corazón. El sacro se observa acortado y el cóccix queda oculto por la angulación caudal. Las articulaciones coxofemorales deben ser visibles. La angulación insuficiente produce un anillo pélvico alargado; el exceso de angulación produce un anillo pélvico cerrado.
		
		\textbf{Ausencia de rotación:} Las alas ilíacas deben ser simétricas en tamaño. Las espinas isquiáticas deben ser simétricas respecto a la línea media (confirman que el RC entró en el plano medio sagital). Pubis e isquion simétricos bilateralmente. Distribución simétrica de todas las estructuras óseas del anillo. La corrección de rotación se realiza nivelando las EIAS.
		
		\textbf{Penetración y contraste:} Detalle trabecular visible en los huesos del anillo pélvico. Diferenciación clara entre las superposiciones óseas normales (ramas del pubis). Visualización de partes blandas adyacentes sin sobreexposición. La técnica intermedia con kVp adecuado debe permitir atravesar la masa del sacro y distinguir estructuras superpuestas.
		
		\subsubsection*{6. Puntos Críticos}
		
		\textbf{Angulación insuficiente:} El anillo pélvico se proyecta alargado en dirección cráneo-caudal, impidiendo la evaluación real del estrecho superior. Corregir aumentando la angulación caudal hasta los 40°.
		
		\textbf{Exceso de angulación:} El anillo pélvico aparece ``cerrado'' o colapsado, con superposición excesiva de estructuras. Reducir la angulación caudal para obtener la apertura correcta del anillo.
		
		\textbf{Adaptaciones:} En pacientes con traumatismo pélvico grave que no toleran el decúbito supino estándar o la flexión de rodillas, se puede realizar en camilla de trauma sin modificar la posición del paciente, ajustando únicamente la angulación del tubo. Si se realiza en decúbito prono, la angulación debe invertirse a \textbf{cefálica} (contraria), ya que en esa posición el sacro se desplaza en lugar de la sínfisis.
		
		\textbf{Trucos:} El RI debe centrarse a la altura de los trocánteres mayores (no en las EIAS) porque el RC entra angulado y se desplaza el punto de proyección. Verificar siempre la simetría de las espinas isquiáticas respecto a la línea media como indicador de centrado correcto en el plano medio sagital. Esta proyección corresponde a una ``vista desde arriba'' de la sínfisis y el anillo pélvico.
		
		
		\onecolumn
		\clearpage
		
		\begin{figure}
			\centering
			\includegraphics[width=0.7\linewidth]{imagenes/pelvis32}
			\caption{}
			\label{fig:pelvis32}
		\end{figure}
		
		
		\twocolumn
		\clearpage
		
		%----------------------------------------------------------
		%  Para agregar otro enfoque duplicá el bloque
		%  desde \subsection hasta \textbf{Trucos}
		%----------------------------------------------------------
	
	
	%%%%%%%%%%%%%%%%%%%%%%%%%%%%%%%%%%%%%%%%%%%%%%%%%%%%%%%%%%%
	%  SECCIÓN: Pelvis
	%  Archivo: secciones/pelvis.tex
	%  Sin preámbulo ni \begin{document} — lo maneja main.tex
		%%%%%%%%%%%%%%%%%%%%%%%%%%%%%%%%%%%%%%%%%%%%%%%%%%%%%%%%%%%
		
		\section{Pelvis}
		
		\subsection{Proyección PA - Pubis / Isquion}
		
		\subsubsection*{1. Identificación}
		
		\textbf{Tipo:} Complementaria.
		
		\textbf{Indicaciones clínicas principales:}
		\begin{itemize}
			\item Traumatismo de la sínfisis pubiana: luxación o diástasis
			\item Patología tumoral o quística del pubis e isquion
			\item Control postoperatorio de reparaciones pélvicas
		\end{itemize}
		
		\subsubsection*{2. Técnica Base}
		
		\textbf{Receptor de imagen:} 24 $\times$ 30 cm. Orientación \textbf{transversal (apaisado)}.
		
		\textbf{DFRI:} 100 cm. Distancia estándar para estudios de pelvis.
		
		\textbf{Técnica:} Intermedia. La posición en decúbito prono permite mayor definición de los bordes óseos, mejor visualización de la parte cortical del hueso y de las estructuras fibrocartilaginosas de la sínfisis del pubis.
		
		\textbf{Bucky:} Sí. Utilizar rejilla antidifusora.
		
		\subsubsection*{3. Posicionamiento}
		
		\paragraph{Paciente}
		Decúbito prono sobre la mesa de examen. El plano medio sagital del cuerpo se centra sobre la línea media de la rejilla. La pelvis no debe presentar rotación.
		
		\begin{figure}[h]
			\centering
			\includegraphics[width=0.7\linewidth]{imagenes/pelvis29}
			\caption{}
			\label{fig:pelvis29}
		\end{figure}
		
		
		\paragraph{Región de interés}
		El receptor se centra a la altura de los trocánteres mayores; esta referencia coincide simultáneamente con la sínfisis del pubis. La región a incluir comprende: mitad inferior del ala ilíaca con las respectivas articulaciones de cadera, borde superior de la sínfisis del pubis, ramas isquiáticas completas y agujeros obturadores en su totalidad.
		
		\paragraph{Rayo Central}
		Perpendicular al plano del receptor. Punto de entrada: \textbf{cóccix distal}; el rayo sale por la sínfisis del pubis.
		
		\paragraph{Receptor de imagen}
		Centrado en el punto medio del receptor, coincidente con la altura de los trocánteres mayores y la sínfisis del pubis. Los bordes del campo deben incluir las alas ilíacas inferiores, las articulaciones coxofemorales y las ramas isquiáticas en su totalidad.
		
		\paragraph{Respiración}
		Suspendida durante la exposición. Evita el movimiento involuntario que podría generar borrosidad en los contornos óseos.
		
		\subsubsection*{4. Anatomía Visible}
		
		Pubis: cuerpo del pubis, rama superior, rama inferior y sínfisis del pubis. Isquión: cuerpo del isquión, rama isquiática, tuberosidad isquiática y espina isquiática. Agujeros obturadores bilaterales, formados por la unión de las ramas del pubis y el isquión. Articulaciones coxofemorales y acetábulos. Sacro y cóccix en relación con la sínfisis púbica. Los huesos pubianos e isquiáticos se visualizan sin magnificación en esta proyección, lo que permite mediciones precisas y evaluación de los bordes de la sínfisis.
		
		\subsubsection*{5. Criterios de Evaluación}
		
		\textbf{Posicionamiento correcto:} Huesos pubianos e isquiáticos próximos al centro de la radiografía y no magnificados. Sacro y cóccix alineados respecto a la sínfisis del pubis. Nitidez en los bordes de la sínfisis, lo que habilita la realización de mediciones. Articulaciones coxofemorales y agujeros obturadores completamente incluidos.
		
		\textbf{Ausencia de rotación:} Agujeros obturadores simétricos bilateralmente. Sacro y cóccix proyectados en la línea media respecto a la sínfisis del pubis.
		
		\textbf{Penetración y contraste:} Visualización adecuada de las ramas pubianas e isquiáticas con contraste suficiente para diferenciar fracturas, diástasis y estructuras fibrocartilaginosas de la sínfisis.
		
		\subsubsection*{6. Puntos Críticos}
		
		\textbf{Rotación del paciente:} Genera asimetría de los agujeros obturadores, imposibilitando la comparación bilateral y enmascarando patología. Corregir alineando el plano medio sagital con la línea media de la rejilla antes de exponer.
		
		\textbf{Centrado incorrecto:} No incluir toda la anatomía de interés --- especialmente las ramas isquiáticas o los agujeros obturadores --- invalida la proyección para diagnóstico de trauma pélvico. El centrado debe realizarse sobre los trocánteres mayores como referencia doble (cadera + sínfisis).
		
		\textbf{Adaptaciones:} En trauma pélvico severo modificar el posicionamiento según las limitaciones del paciente, priorizando la inmovilidad de fracturas. En pacientes con dolor pélvico utilizar soportes bajo el abdomen sin comprometer la alineación. En embarazo evaluar necesidad clínica versus riesgo de radiación fetal.
		
		\textbf{Trucos:} La proyección PA es complementaria y su valor diferencial frente al AP axial de Taylor es que los huesos no están magnificados, permitiendo mediciones fiables de la diástasis. La nitidez del borde de la sínfisis es el parámetro clave cuando se sospecha diástasis. Para trauma con desplazamiento vertical considerar las proyecciones outlet (desplazamiento supero-inferior) e inlet (desplazamiento posterior) como complementos adicionales.
		
	
	\onecolumn
	\clearpage
	
	\begin{figure}
		\centering
		\includegraphics[width=0.7\linewidth]{imagenes/pelvis30}
		\caption{}
		\label{fig:pelvis30}
	\end{figure}
	
	
	\twocolumn
	\clearpage
		
		%----------------------------------------------------------
		%  Para agregar otro enfoque duplicá el bloque
		%  desde \subsection hasta \textbf{Trucos}
		%----------------------------------------------------------

%%%%%%%%%%%%%%%%%%%%%%%%%%%%%%%%%%%%%%%%%%%%%%%%%%%%%%%%%%%
%  SECCIÓN: Pelvis
%  Archivo: secciones/pelvis.tex
%  Sin preámbulo ni \begin{document} — lo maneja main.tex
	%%%%%%%%%%%%%%%%%%%%%%%%%%%%%%%%%%%%%%%%%%%%%%%%%%%%%%%%%%%
	
	\section{Pelvis}
	
	\subsection{Proyección Oblicua Alar - Método de Judet}
	
	\subsubsection*{1. Identificación}
	
	\textbf{Tipo:} Proyección complementaria. Forma par con la proyección oblicua obturatriz (Judet interno).
	
	\textbf{Indicaciones clínicas principales:}
	\begin{itemize}
		\item Fracturas del ala ilíaca y columna posterior del acetábulo (línea ilioisquiática)
		\item Traumatismo pélvico con sospecha de compromiso acetabular
		\item Evaluación del borde anterior del acetábulo
		\item Visualización de escotaduras ciáticas mayor y menor
		\item Evaluación de la articulación sacroilíaca
		\item Acetábulo en perfil
	\end{itemize}
	
	\subsubsection*{2. Técnica Base}
	
	\textbf{Receptor de imagen:} 24 $\times$ 30 cm. Orientación \textbf{longitudinal}. Borde superior por encima de la cresta ilíaca del lado afectado.
	
	\textbf{DFRI:} 100 cm.
	
	\textbf{Técnica:} Intermedia. Permite visualizar la cabeza femoral a través del hueso acetabular, diferenciar trabeculado óseo y cortical, y evidenciar las líneas ilioisquiáticas en detalle.
	
	\textbf{Bucky:} Bucky de mesa con grilla.
	
	\subsubsection*{3. Posicionamiento}
	
	\paragraph{Paciente}
	Decúbito semisupino apoyado sobre el lado afectado. El plano sagital del paciente se centra de modo que la cadera del lado en estudio quede alineada con la línea media del receptor.
	
	\begin{figure}[h]
		\centering
		\includegraphics[width=0.7\linewidth]{imagenes/pelvis33}
		\caption{}
		\label{fig:pelvis33}
	\end{figure}
	
	
	\paragraph{Región de interés}
	Se realiza una rotación de 45° respecto al plano de la mesa, elevando el lado no afectado (sano). El objetivo es desplegar el ala ilíaca del lado afectado paralela al plano del receptor (OPD u OPI según el lado). Se apoya el hombro, cadera y rodilla elevados sobre sacos de arena para estabilizar la posición. Se ajustan las crestas ilíacas al mismo plano transverso verificando la EIAS como referencia.
	
	\paragraph{Rayo Central}
	Perpendicular al plano del receptor. Punto de entrada: \textbf{5 cm distal y 5 cm medial a la espina ilíaca anterosuperior (EIAS) del lado afectado}.
	
	\paragraph{Receptor de imagen}
	Centrado a la altura de la EIAS. Se incluye desde el ala ilíaca afectada hasta la región proximal del fémur; lateralmente desde la sínfisis púbica hasta el borde cutáneo lateral, con la articulación coxofemoral centrada en el campo.
	
	\paragraph{Respiración}
	Suspendida en espiración. Elimina el movimiento involuntario y evita borrosidad de estructuras óseas pelvianas finas.
	
	\subsubsection*{4. Anatomía Visible}
	
	Se visualiza el ala ilíaca desplegada y ancha sin superposiciones, con las escotaduras ciáticas mayor y menor sin obstrucciones. Se proyecta el acetábulo en perfil, la columna ilioisquiática posterior y el borde anterior del acetábulo. Son visibles también la articulación sacroilíaca, la cabeza femoral dentro del acetábulo y la región craneal del ala del ilíaco. El foramen obturador aparece cerrado o perfilado, lo que confirma la rotación correcta de 45°.
	
	\subsubsection*{5. Criterios de Evaluación}
	
	\textbf{Posicionamiento correcto:} Ala ilíaca del lado afectado se proyecta ancha y desplegada; foramen obturador cerrado; columna posterior en perfil; articulación sacroilíaca y coxofemoral visibles sin superposición mayor.
	
	\textbf{Ausencia de rotación:} El ala ilíaca no debe verse estrecha (indicaría rotación insuficiente); las escotaduras ciáticas deben estar libres de superposición; el acetábulo se aprecia en perfil verdadero.
	
	\textbf{Penetración y contraste:} Contraste intermedio que permita diferenciar trabeculado óseo, cortical y cabeza femoral a través del hueso acetabular; escala de grises suficiente para evidenciar líneas de fractura en la columna ilioisquiática y en el borde anterior del acetábulo.
	
	\subsubsection*{6. Puntos Críticos}
	
	\textbf{Rotación insuficiente:} El ala ilíaca del lado apoyado se ve estrecha y el foramen obturador no queda cerrado, impidiendo el perfil real del acetábulo. Corregir elevando más el lado sano hasta alcanzar los 45°.
	
	\textbf{Centrado incorrecto del RC:} Si el RC no se dirige 5 cm distal y 5 cm medial a la EIAS, se pierde la tangencia sobre la columna posterior y el borde anterior del acetábulo queda distorsionado. Verificar la referencia anatómica antes de exponer.
	
	\textbf{Adaptaciones:} En pacientes con trauma agudo y limitación para la rotación activa, utilizar sacos de arena o cuñas para sostener la posición oblicua; reducir el ángulo si el dolor impide los 45° completos y documentar el ángulo real obtenido.
	
	\textbf{Trucos:} Confirmar la correcta rotación observando el ala ilíaca antes de exponer: si se ve ancha y el agujero obturador está cerrado, la posición es válida. Esta proyección se denomina \textit{alar} porque lo que se despliega es precisamente el ala del ilion del lado afectado; la rotación siempre es hacia ese mismo lado.
	
\onecolumn
\clearpage

\begin{figure}
	\centering
	\includegraphics[width=0.7\linewidth]{imagenes/pelvis35}
	\caption{}
	\label{fig:pelvis35}
\end{figure}


\twocolumn
\clearpage
	
	%----------------------------------------------------------
	%  Para agregar otro enfoque duplicá el bloque
	%  desde \subsection hasta \textbf{Trucos}
	%----------------------------------------------------------
	
	
%%%%%%%%%%%%%%%%%%%%%%%%%%%%%%%%%%%%%%%%%%%%%%%%%%%%%%%%%%%
%  SECCIÓN: Pelvis — Proyecciones de Judet
%  Archivo: secciones/judet_oblicua_obturatriz.tex
%  Sin preámbulo ni \begin{document} — lo maneja main.tex
	%%%%%%%%%%%%%%%%%%%%%%%%%%%%%%%%%%%%%%%%%%%%%%%%%%%%%%%%%%%
	
	\section{Pelvis}
	
	\subsection{Oblicua Interna u Obturatriz (Judet)}
	
	\subsubsection*{1. Identificación}
	
	\textbf{Tipo:} Proyección complementaria. Forma par con la Oblicua Externa o Alar (Judet).
	
	\textbf{Indicaciones clínicas principales:}
	\begin{itemize}
		\item Fracturas del acetábulo: columna anterior (iliopubiana) y pared posterior
		\item Trauma pélvico complejo con sospecha de compromiso acetabular
		\item Fracturas de la rama isquiopubiana (borde inferior del agujero obturador)
		\item Evaluación del borde posterior del acetábulo y la columna iliopubiana
	\end{itemize}
	
	\subsubsection*{2. Técnica Base}
	
	\textbf{Receptor de imagen:} 24$\times$30 cm. Orientación \textbf{longitudinal} en paciente estándar.
	
	\textbf{DFRI:} 100 cm. Distancia estándar para pelvis con Bucky mural o de mesa.
	
	\textbf{Técnica:} Intermedia. Permite visualizar la cabeza femoral a través del hueso acetabular y evaluar las líneas iliopubiana e ilioisquiática por transparencia; una técnica dura impediría apreciar si la línea cortical está interrumpida, y una blanda no penetraría suficientemente el macizo óseo.
	
	\textbf{Bucky:} Sí, con grilla.
	
	\subsubsection*{3. Posicionamiento}
	
	\paragraph{Paciente}
	Decúbito semisupino con el \textbf{lado afectado elevado}. El cuerpo se apoya sobre el lado sano; la cadera en estudio queda rotada hacia arriba, alejada de la mesa. A diferencia de la oblicua alar, en esta proyección se radiografía el lado \textbf{elevado} (no el apoyado).
	
	\begin{figure}[h]
		\centering
		\includegraphics[width=0.7\linewidth]{imagenes/pelvis36}
		\caption{}
		\label{fig:pelvis36}
	\end{figure}
	
	
	\paragraph{Región de interés}
	Se eleva el lado afectado hasta que la superficie anterior del cuerpo forme un \textbf{ángulo de 45°} con la mesa. La cadera en estudio se centra sobre la línea media del receptor. Se estabiliza la posición con soportes apropiados bajo el hemitórax y la rodilla elevada.
	
	\paragraph{Rayo Central}
	Perpendicular al plano del receptor. Punto de entrada: \textbf{5 cm distal a la EIAS del lado afectado elevado}.
	
	\paragraph{Receptor de imagen}
	Centrado en la articulación coxofemoral del lado afectado. La región a incluir abarca desde el ala ilíaca afectada hasta la porción proximal del fémur; lateralmente, desde la sínfisis púbica hasta el borde cutáneo ipsilateral.
	
	\paragraph{Respiración}
	Suspendida en espiración. Elimina el movimiento involuntario y evita artefactos cinéticos en la imagen pélvica.
	
	\subsubsection*{4. Anatomía Visible}
	
	En la oblicua obturatriz el lado elevado queda orientado de frente al haz, por lo que el \textbf{agujero obturador del lado afectado} se proyecta en su \textbf{máxima apertura}, permitiendo evaluar la integridad de la rama isquiopubiana en toda su extensión. El \textbf{ala ilíaca del mismo lado} se proyecta de perfil, estrecha y superpuesta sobre sí misma. Se visualizan con claridad la \textbf{columna iliopubiana (anterior)} —que se extiende desde la sínfisis del pubis y el agujero obturador hasta la EIAS— y el \textbf{borde posterior del acetábulo}. La columna ilioisquiática (posterior) también es visible: compuesta por la porción vertical del isquion y la parte del ilion inmediatamente superior. Se observan además la \textbf{cabeza femoral} dentro del acetábulo y la \textbf{articulación coxofemoral} centradas en el receptor.
	
	\subsubsection*{5. Criterios de Evaluación}
	
	\textbf{Posicionamiento correcto:} El agujero obturador del lado afectado (elevado) aparece abierto en su máxima apertura, sin superposición de las ramas pubianas e isquiáticas entre sí. El ala ilíaca del lado elevado se proyecta de perfil, estrecha. La cabeza femoral debe estar dentro del acetábulo y las espinas ilíacas superpuestas confirman la rotación adecuada de 45°.
	
	\textbf{Ausencia de rotación:} Rotación correcta equivale a agujero obturador completamente abierto. Si el agujero aparece cerrado, hay exceso de rotación (más de 45°). Si el ala ilíaca afectada se ve ancha y desplegada, falta rotación (menos de 45°, imagen similar a la oblicua alar).
	
	\textbf{Penetración y contraste:} Técnica intermedia que permita ver la cabeza femoral a través del acetábulo y que las líneas corticales iliopubiana e ilioisquiática sean visibles por transparencia. Una técnica insuficiente vuelve la placa radiopaca e impide evaluar la integridad cortical; una técnica excesiva quema las estructuras superficiales.
	
	\subsubsection*{6. Puntos Críticos}
	
	\textbf{Exceso de rotación (más de 45°):} El agujero obturador aparece cerrado o semicerrado, impidiendo evaluar la rama isquiopubiana. Corrección: reducir el ángulo de elevación del lado afectado.
	
	\textbf{Falta de rotación (menos de 45°):} El ala ilíaca del lado afectado se ve ancha y desplegada; la imagen se asemeja a la oblicua alar, perdiendo la apertura obturatriz. Corrección: aumentar la elevación hasta los 45° exactos.
	
	\textbf{Adaptaciones:} En pacientes politraumatizados que no toleran la rotación activa, el ángulo puede lograrse con cuñas radiotransparentes. Si existe inestabilidad pélvica severa, la maniobra debe coordinarse con el equipo clínico antes de posicionar.
	
	\textbf{Trucos:} Recordar la regla práctica: en Judet se radiografía el lado \textbf{elevado} en la obturatriz y el lado \textbf{apoyado} en la alar. El dato diagnóstico clave de esta proyección es el agujero obturador abierto; si no está abierto, la proyección no cumple su propósito. La rama isquiopubiana corresponde al borde inferior del agujero obturador — dato que suele confundirse en el oral.
	
	\onecolumn
	\clearpage
	
	\begin{figure}
		\centering
		\includegraphics[width=0.7\linewidth]{imagenes/pelvis37}
		\caption{}
		\label{fig:pelvis37}
	\end{figure}
	
	\twocolumn
	\clearpage
	
	%----------------------------------------------------------
	%  Para agregar otro enfoque duplicá el bloque
	%  desde \subsection hasta \textbf{Trucos}
	%----------------------------------------------------------
	
	
%%%%%%%%%%%%%%%%%%%%%%%%%%%%%%%%%%%%%%%%%%%%%%%%%%%%%%%%%%%
%  SECCIÓN: Pelvis — Cadera
%  Archivo: secciones/pelvis_cadera.tex
%  Sin preámbulo ni \begin{document} — lo maneja main.tex
	%%%%%%%%%%%%%%%%%%%%%%%%%%%%%%%%%%%%%%%%%%%%%%%%%%%%%%%%%%%
	
	\section{Pelvis}
	
	\subsection{Enfoque de Dunn 45°}
	
	\subsubsection*{1. Identificación}
	
	\textbf{Tipo:} Proyección complementaria. Forma par con el enfoque de Dunn 90° y la radiografía de frente AP de pelvis.
	
	\textbf{Indicaciones clínicas principales:}
	\begin{itemize}
		\item Pinzamiento femoroacetabular (FAI): unilateral o bilateral comparativa
		\item Dolor inguinal crónico en pacientes jóvenes o deportistas
		\item Evaluación de cobertura acetabular y transición cabeza-cuello femoral
		\item Displasia de cadera y deformidad tipo CAM o PINCER
		\item Osteonecrosis de cabeza femoral y valoración del espacio articular coxofemoral
	\end{itemize}
	
	\subsubsection*{2. Técnica Base}
	
	\textbf{Receptor de imagen:} 24$\times$30 cm (estudio unilateral) o 35$\times$43 cm (estudio bilateral). Orientación \textbf{transversal} en ambos casos.
	
	\textbf{DFRI:} 100 cm. Distancia estándar para estudios de pelvis y cadera.
	
	\textbf{Técnica:} Intermedia. Permite diferenciar el trabeculado óseo y las corticales para detectar detalles finos en la unión cabeza-cuello femoral.
	
	\textbf{Bucky:} Bucky con grilla. Se emplea para reducir la radiación dispersa en una región de gran volumen tisular.
	
	\subsubsection*{3. Posicionamiento}
	
	\paragraph{Paciente}
	Decúbito supino sobre la mesa de examen. En el estudio bilateral, el plano medio sagital del cuerpo se centra en la línea media de la rejilla.
	
	\begin{figure}[h]
		\centering
		\includegraphics[width=0.7\linewidth]{imagenes/pelvis38}
		\caption{}
		\label{fig:pelvis38}
	\end{figure}
	
	
	\paragraph{Región de interés}
	Cadera y rodillas en flexión de 45°, lo que revela las alteraciones de la porción anterior de la unión cabeza-cuello femoral. Se aplica una abducción del muslo de 20°: las rodillas se separan entre sí y se bajan hacia afuera. Los pies quedan apoyados en posición neutra, con la tibia paralela al eje longitudinal del cuerpo.
	
	\paragraph{Rayo Central}
	Perpendicular al plano del receptor.
	\begin{itemize}
		\item \textbf{Unilateral:} RC perpendicular al cuello femoral del lado estudiado.
		\item \textbf{Bilateral:} RC perpendicular centrado en el punto medio entre la sínfisis del pubis y la EIAS.
	\end{itemize}
	
	\paragraph{Receptor de imagen}
	Centrado al mismo nivel que la radiografía de frente AP de pelvis.
	\begin{itemize}
		\item \textbf{Unilateral:} incluir mitad inferior del ala ilíaca y borde superior del acetábulo; inferiormente, el tercio proximal del fémur; medialmente, la sínfisis del pubis; lateralmente, las partes blandas del muslo afectado.
		\item \textbf{Bilateral:} incluir ambas EIAS superiormente; inferiormente, el tercio proximal del fémur y parte de las diáfisis proximales; lateralmente, ambos bordes cutáneos.
	\end{itemize}
	
	\paragraph{Respiración}
	Suspendida durante la exposición.
	
	\subsubsection*{4. Anatomía Visible}
	
	Se despliega la unión cabeza-cuello anterosuperior del fémur. El cuello femoral se visualiza alargado cuando la flexión de cadera es correcta. Se aprecia el contorno de la cabeza femoral, la unión cabeza-cuello, el borde acetabular anterior y el labrum acetabular de forma indirecta. Son evaluables el offset cabeza-cuello femoral, el ángulo alfa ($\alpha$), la cobertura acetabular y el espacio articular coxofemoral. El trocánter mayor queda por fuera y lateral al cuello; el trocánter menor se demuestra de perfil cuando la abducción es correcta y el pie permanece neutro. La Dunn 45° proporciona una visualización intermedia de la unión cabeza-cuello, a diferencia de la Dunn 90° que maximiza la visualización del contorno anterior.
	
	\subsubsection*{5. Criterios de Evaluación}
	
	\textbf{Posicionamiento correcto:} Cuello femoral alargado; trocánter menor demostrado de perfil; trocánter mayor lateral y separado del cuello; unión cabeza-cuello anterosuperior desplegada; cabeza femoral centrada en el acetábulo; espacio articular coxofemoral uniforme y simétrico; contorno de la cabeza femoral claramente definido.
	
	\textbf{Ausencia de rotación:} En el estudio bilateral, simetría entre ambas caderas. Trocánter menor visible de perfil confirma abducción de 20° y pie en posición neutra.
	
	\textbf{Penetración y contraste:} Técnica intermedia que permite diferenciar el trabeculado óseo y las corticales; densidad suficiente para detectar detalles finos en la transición cabeza-cuello.
	
	\subsubsection*{6. Puntos Críticos}
	
	\textbf{Flexión de cadera insuficiente:} El cuello femoral aparece corto; la unión cabeza-cuello no se despliega correctamente. Aumentar la flexión hasta los 45° verificando que la pierna no quede demasiado plana.
	
	\textbf{Exceso de abducción:} El trocánter mayor superpone el cuello femoral, impidiendo su evaluación. Reducir la abducción hasta los 20° llevando el muslo hacia la línea media.
	
	\textbf{Adaptaciones:} En pacientes con limitación de movilidad o dolor agudo, ajustar el grado de flexión y abducción dentro de la tolerancia; documentar el ángulo real obtenido. Se puede realizar como alternativa a la posición de cross table.
	
	\textbf{Trucos:} Verificar que el pie quede en posición neutra: la rotación interna hace desaparecer el trocánter menor y no permite confirmar la abducción correcta. Confirmar que el trocánter mayor quede lateral y separado del cuello antes de exponer. El centrado es idéntico al de la AP de pelvis de frente, lo que facilita la reproducibilidad.
	
	\clearpage
	\onecolumn
	
	\begin{figure}
		\centering
		\includegraphics[width=0.7\linewidth]{imagenes/pelvis40}
		\caption{Dunn a 45 grados}
		\label{fig:pelvis40}
	\end{figure}
	
	
	\twocolumn
	\clearpage
	
	%----------------------------------------------------------
	
	\subsection{Enfoque de Dunn 90°}
	
	\subsubsection*{1. Identificación}
	
	\textbf{Tipo:} Proyección complementaria. Forma par con el enfoque de Dunn 45° y la radiografía de frente AP de pelvis.
	
	\textbf{Indicaciones clínicas principales:}
	\begin{itemize}
		\item Máxima visualización del contorno anterior de la cabeza femoral
		\item Displasia de cadera y pinzamiento femoroacetabular (FAI)
		\item Deformidad tipo CAM o PINCER
		\item Osteonecrosis de cabeza femoral
		\item Valoración del espacio articular coxofemoral
	\end{itemize}
	
	\subsubsection*{2. Técnica Base}
	
	\textbf{Receptor de imagen:} 24$\times$30 cm (estudio unilateral) o 35$\times$43 cm (estudio bilateral). Orientación \textbf{transversal} en ambos casos.
	
	\textbf{DFRI:} 100 cm. Distancia estándar para estudios de pelvis y cadera.
	
	\textbf{Técnica:} Intermedia. Acorde al biotipo del paciente. Permite diferenciar el trabeculado óseo y las corticales.
	
	\textbf{Bucky:} Bucky con grilla.
	
	\subsubsection*{3. Posicionamiento}
	
	\paragraph{Paciente}
	Decúbito supino sobre la mesa de examen. El plano medio sagital del cuerpo se centra en la línea media de la rejilla.
	
	\begin{figure}[h]
		\centering
		\includegraphics[width=0.7\linewidth]{imagenes/pelvis39}
		\caption{}
		\label{fig:pelvis39}
	\end{figure}
	
	
	\paragraph{Región de interés}
	Caderas flexionadas a 90° (fémures en posición vertical). Se mantiene una abducción de 20°. La tibia debe quedar paralela al eje longitudinal del cuerpo. Puede realizarse de forma unilateral o bilateral.
	
	\paragraph{Rayo Central}
	Perpendicular al receptor. Centrado en el punto medio entre la sínfisis del pubis y la EIAS.
	
	\paragraph{Receptor de imagen}
	Centrado al mismo nivel que la radiografía de frente AP de pelvis, siguiendo el mismo punto de referencia descrito para el RC.
	
	\paragraph{Respiración}
	Suspendida durante la exposición.
	
	\subsubsection*{4. Anatomía Visible}
	
	La Dunn 90° proporciona la máxima visualización del contorno anterior de la cabeza femoral. Se aprecia el offset cabeza-cuello completamente evaluable, la deformidad tipo CAM claramente demostrada si está presente, el borde acetabular anterior, el labrum acetabular de forma indirecta y el espacio articular coxofemoral. Son medibles el ángulo alfa ($\alpha$) y la cobertura acetabular. A mayor flexión respecto de la Dunn 45°, mayor exposición de la zona anterior de la unión cabeza-cuello femoral.
	
	\subsubsection*{5. Criterios de Evaluación}
	
	\textbf{Posicionamiento correcto:} Fémures verticales (flexión de 90°); abducción de 20° mantenida; tibia paralela al eje del cuerpo; contorno anterior de la cabeza femoral completamente visible; unión cabeza-cuello sin superposiciones; offset cabeza-cuello completamente evaluable.
	
	\textbf{Ausencia de rotación:} En el estudio bilateral, simetría entre ambas caderas. Cabeza femoral centrada en el acetábulo; espacio articular uniforme y simétrico.
	
	\textbf{Penetración y contraste:} Técnica intermedia acorde al biotipo; densidad adecuada para demostrar deformidad tipo CAM si está presente y para evaluar el contorno óseo con detalle fino.
	
	\subsubsection*{6. Puntos Críticos}
	
	\textbf{Flexión insuficiente (menor de 90°):} No se obtiene la máxima exposición del contorno anterior de la cabeza femoral; la imagen se asemeja a la Dunn 45° sin aportar información adicional. Verificar que los fémures queden verticales.
	
	\textbf{Abducción incorrecta:} Una abducción mayor de 20° superpone el trocánter mayor sobre el cuello; una abducción insuficiente puede generar superposición de estructuras mediales. Mantener los 20° indicados.
	
	\textbf{Adaptaciones:} En pacientes con limitación de movilidad de cadera, documentar el ángulo de flexión real alcanzado. Ajustar la técnica según biotipo del paciente.
	
	\textbf{Trucos:} Comparar siempre con la Dunn 45° del mismo paciente: la diferencia entre ambos ángulos de flexión permite evaluar la progresión de la deformidad en distintos rangos de movimiento. Verificar la posición de la tibia antes de exponer.
	
	\onecolumn
	\clearpage
	
	\begin{figure}
		\centering
		\includegraphics[width=0.7\linewidth]{imagenes/pelvis41}
		\caption{}
		\label{fig:pelvis41}
	\end{figure}
	
	
	\twocolumn
	\clearpage
	
	%----------------------------------------------------------
	%  Para agregar otro enfoque duplicá el bloque
	%  desde \subsection hasta \textbf{Trucos}
	%----------------------------------------------------------
	
	
%%%%%%%%%%%%%%%%%%%%%%%%%%%%%%%%%%%%%%%%%%%%%%%%%%%%%%%%%%%
%  SECCIÓN: Pelvis — Posición de Von Rosen
%  Archivo: secciones/von_rosen.tex
%  Sin preámbulo ni \begin{document} — lo maneja main.tex
	%%%%%%%%%%%%%%%%%%%%%%%%%%%%%%%%%%%%%%%%%%%%%%%%%%%%%%%%%%%
	
	\section{Pelvis}
	
	\subsection{Posición de Von Rosen (Cadera Pediátrica)}
	
	\subsubsection*{1. Identificación}
	
	\textbf{Tipo:} Complementario especializado. Técnica de posicionamiento forzado para evaluación de la cadera en pacientes pediátricos.
	
	\textbf{Indicaciones clínicas principales:}
	\begin{itemize}
		\item Sospecha de luxación congénita o displasia del desarrollo de cadera (DDH)
		\item Inestabilidad de cadera en neonatos y lactantes
		\item Evaluación de reducción de cadera luxada (verificar que la cabeza femoral vuelve a su lugar)
		\item Asimetría de pliegues o acortamiento aparente de miembros inferiores
	\end{itemize}
	
	\subsubsection*{2. Técnica Base}
	
	\textbf{Receptor de imagen:} 24$\times$30 cm, orientación \textbf{longitudinal}. El tamaño puede variar según el tamaño del paciente.
	
	\textbf{DFRI:} 100 cm. Técnica estándar para proyecciones de pelvis pediátrica.
	
	\textbf{Técnica:} Homogeneizada (RP). Se utilizan valores mínimos para distinguir tejido óseo y partes blandas, permitiendo visualizar los núcleos de osificación (en su mayoría cartilaginosos en neonatos) sin sobreexponer.
	
	\textbf{Bucky:} Sin Bucky ni rejilla. Se evita para reducir la dosis de radiación, dado que se trata de gónadas en desarrollo.
	
	\subsubsection*{3. Posicionamiento}
	
	\paragraph{Paciente}
	Decúbito supino sobre la mesa de examen. Sin pañal. Se recomienda ambiente cálido para evitar llanto y movimiento. Para el posicionamiento bilateral simétrico se requieren al menos dos personas.
	
	\begin{figure}[h]
		\centering
		\includegraphics[width=0.7\linewidth]{imagenes/pelvis42}
		\caption{}
		\label{fig:pelvis42}
	\end{figure}
	
	
	\paragraph{Región de interés}
	Ambas caderas en abducción forzada a un mínimo de 45° (óptimo entre 65° y 75°). Los fémures se rotan internamente al máximo hasta que las plantas de los pies se apoyan entre sí. Caderas y rodillas en extensión completa. El posicionamiento debe ser bilateral y simétrico. La abducción a 45° hace que el cuello femoral quede paralelo al receptor de imagen, de modo que el fémur se visualiza de frente.
	
	\paragraph{Rayo Central}
	Perpendicular al plano del receptor. Punto de entrada: \textbf{2 dedos por encima de la sínfisis del pubis}, centrado en la línea media entre la sínfisis del pubis y la EIAS.
	
	\paragraph{Receptor de imagen}
	Centrado en la región de la sínfisis del pubis y caderas bilaterales. Debe incluir desde la cresta ilíaca hasta las diáfisis femorales proximales, para permitir las mediciones de líneas y ángulos diagnósticos.
	
	\paragraph{Respiración}
	No aplicable. Se trata de pacientes neonatos o lactantes. La exposición debe realizarse durante un período de inmovilidad, evitando exponer durante el llanto o la agitación del paciente.
	
	\subsubsection*{4. Anatomía Visible}
	
	Se visualizan las articulaciones coxofemorales bilaterales, las cabezas femorales y su posición relativa al acetábulo, los núcleos de osificación femoral cuando están presentes, las alas ilíacas, los agujeros obturados, las articulaciones sacroilíacas y la sínfisis del pubis. Se evalúa la relación entre el eje longitudinal del fémur y el acetábulo, la cobertura acetabular y la posición de la cabeza femoral respecto a los cuadrantes de Hilgenreiner--Perkins. En un paciente con cadera luxada, el fémur afectado se desplaza hacia arriba y hacia afuera; la abducción forzada hace más evidente este desplazamiento.
	
	\subsubsection*{5. Criterios de Evaluación}
	
	\textbf{Posicionamiento correcto:} Simetría de ambas alas ilíacas y agujeros obturados (pelvis sin rotación). Ambos fémures con aspecto frontal producto de la abducción mínima de 45°. Inclusión completa desde crestas ilíacas hasta diáfisis femorales proximales.
	
	\textbf{Ausencia de rotación:} Simetría bilateral de la pelvis ósea. Igual grado de abducción en ambos miembros. Si hay asimetría entre ambos lados, sospechar error de posicionamiento antes de interpretar como patología.
	
	\textbf{Penetración y contraste:} Visualización nítida de las estructuras óseas y de los núcleos de osificación. Contraste suficiente para distinguir partes blandas. Sin sobreexposición que oculte detalles finos del cartílago o los núcleos de osificación.
	
	\subsubsection*{6. Puntos Críticos}
	
	\textbf{Abducción insuficiente:} Si la abducción es menor de 45°, el cuello femoral no queda paralelo al receptor y pueden generarse falsos positivos. Verificar visualmente el ángulo antes de exponer.
	
	\textbf{Asimetría de posición:} Diferente grado de abducción o rotación entre ambos miembros invalida la comparación bilateral. Corregir con asistencia de una segunda persona antes de exponer.
	
	\textbf{Adaptaciones:} Para pacientes agitados, realizar la exposición durante el sueño o con el bebé calmado. Utilizar dispositivos de posicionamiento para mantener los ángulos correctos sin fuerza excesiva.
	
	\textbf{Trucos:} Calentar previamente el ambiente para relajar al paciente. Confirmar simetría visual antes de la exposición. Aplicar protección gonadal cuando sea posible sin ocultar la anatomía de interés. La interpretación debe correlacionarse siempre con los hallazgos clínicos y, cuando esté disponible, con estudios ecográficos (preferidos en menores de 4 meses).
	
	\onecolumn
	\clearpage
	
	\begin{figure}
		\centering
		\includegraphics[width=0.7\linewidth]{imagenes/pelvis45}
		\caption{}
		\label{fig:pelvis45}
	\end{figure}
	
	\twocolumn
	\clearpage
	
	\subsubsection*{Información Adicional}
	
	\paragraph{Principio diagnóstico de la técnica}
	La posición de Von Rosen provoca la luxación de la(s) cadera(s) comprometida(s) mediante la abducción forzada a 45°, reproduciendo el equivalente radiográfico del test de Ortolani. Las caderas con subluxación se reducen en esta posición, mientras que las completamente luxadas no lo hacen.
	
	\paragraph{Interpretación del eje femoral}
	En un niño normal, el eje longitudinal del fémur se alinea con el margen lateral del acetábulo y la cabeza femoral se encuentra en el cuadrante infero-interno (cuadrantes de Hilgenreiner--Perkins). En un niño con cadera displásica o luxada, el eje longitudinal del fémur se orienta hacia la espina ilíaca anterosuperior (EIAS).
	
	\paragraph{Líneas diagnósticas}
	Se utilizan las líneas de Hilgenreiner (horizontal, a través de ambos cartílagos trirradiados) y de Perkins (vertical, desde el borde lateral del acetábulo), que dividen cada cadera en cuatro cuadrantes. El núcleo de osificación de la cabeza femoral debe ubicarse en el cuadrante infero-interno. En luxación, se desplaza al cuadrante supero-externo.
	
	\paragraph{Ángulo acetabular}
	Mide la inclinación del techo acetabular. El valor normal al nacer es de 30°. Se considera patológico si supera los 36° al nacer o los 30° a los 3 meses de edad. Disminuye entre 0{,}5° y 1° por mes; debe ser aproximadamente 20° a los 2 años, reflejando la maduración progresiva del acetábulo.
	
	\paragraph{Consideraciones de edad y modalidad}
	Esta técnica es óptima entre los 0 y 6 meses de vida. A partir de los 8 meses la osificación avanzada dificulta la interpretación. En menores de 4 meses se prefiere la ecografía como primera línea diagnóstica, dado que gran parte de la pelvis es cartilaginosa y la exposición a radiación debe minimizarse.
	
	\begin{figure}[h]
		\centering
		\includegraphics[width=0.7\linewidth]{imagenes/pelvis44}
		\caption{}
		\label{fig:pelvis44}
	\end{figure}
	
	\begin{figure}[h]
		\centering
		\includegraphics[width=0.7\linewidth]{imagenes/pelvis43}
		\caption{}
		\label{fig:pelvis43}
	\end{figure}
	
	\clearpage
	
	%----------------------------------------------------------
	%  Para agregar otro enfoque duplicá el bloque
	%  desde \subsection hasta \textbf{Trucos}
	%----------------------------------------------------------
	
	
%%%%%%%%%%%%%%%%%%%%%%%%%%%%%%%%%%%%%%%%%%%%%%%%%%%%%%%%%%%
%  SECCIÓN: Pelvis
%  Archivo: secciones/pelvis.tex
%  Bloque adicional — Signo de Risser
%  Sin preámbulo ni \begin{document} — lo maneja main.tex
	%%%%%%%%%%%%%%%%%%%%%%%%%%%%%%%%%%%%%%%%%%%%%%%%%%%%%%%%%%%
	\section{Pelvis}
	
	\subsection{Signo de Risser}
	
	\subsubsection*{Definición y Propósito}
	
	El \textbf{Signo de Risser} es un indicador radiográfico utilizado para evaluar la madurez
	esquelética mediante el análisis del grado de osificación de la \textbf{apófisis ilíaca
		(cresta ilíaca)} en radiografías de pelvis. Fue desarrollado por el Dr.\ Joseph C.\ Risser
	en 1958.
	
	\textbf{Objetivo principal:} predecir el potencial de crecimiento vertebral y la
	progresión de escoliosis en adolescentes.
	
	\textbf{Proyección requerida:} AP de pelvis, preferiblemente en bipedestación, con
	visualización completa de ambas crestas ilíacas. Habitualmente incluida en estudios
	de columna completa.
	
	\subsubsection*{Fundamento Anatómico}
	
	La osificación de la apófisis ilíaca se inicia en la región \textbf{anterolateral} de
	la cresta ilíaca (cerca de la espina ilíaca anterosuperior), avanza medialmente hacia
	la espina ilíaca posterosuperior y se fusiona posteriormente en dirección
	\textbf{medial a lateral}.
	
	Esta osificación ocurre de manera paralela a la osificación de las apófisis
	vertebrales, lo que permite estimar el potencial de crecimiento espinal restante.
	
	\subsubsection*{Sistema de Clasificación — Escala Risser (0--5)}
	
	\textbf{Risser 0:} sin osificación visible de la apófisis ilíaca. Indica crecimiento
	significativo restante. Riesgo alto de progresión de escoliosis. Etapa prepuberal o
	pubertad temprana.
	
	\textbf{Risser 1:} osificación del 25\% de la cresta ilíaca. Etapa prepuberal o
	pubertad temprana. Considerable potencial de crecimiento restante.
	
	\textbf{Risser 2:} osificación del 50\% de la cresta ilíaca. Antes o durante el pico
	de crecimiento. Moderado potencial de crecimiento restante.
	
	\textbf{Risser 3:} osificación del 75\% de la cresta ilíaca. Etapa de desaceleración
	del crecimiento. Limitado potencial de crecimiento restante.
	
	\textbf{Risser 4:} osificación completa (100\%) sin fusión. Casi cesación del
	crecimiento. Mínimo potencial de crecimiento restante.
	
	\textbf{Risser 5:} fusión completa de la apófisis con el ilion. Madurez esquelética
	completa. No hay crecimiento restante.
	
	\subsubsection*{Criterios de Evaluación Radiográfica}
	
	\textbf{Identificación:} visualización clara y bilateral de la apófisis ilíaca en
	ambas crestas.
	
	\textbf{Medición:} estimación del porcentaje de osificación respecto a la longitud
	total de la cresta ilíaca.
	
	\textbf{Fusión:} detección de unión apofisaria completa al ilion (estadio 5).
	
	\subsubsection*{Aplicaciones Clínicas}
	
	\textbf{Evaluación de escoliosis:} los pacientes con Risser 0--2 presentan mayor riesgo
	de progresión. El estadio determina la urgencia de intervención y orienta el
	seguimiento durante el crecimiento.
	
	\textbf{Decisiones terapéuticas:}
	\begin{itemize}
		\item \textbf{Corsé ortopédico:} más efectivo en grados Risser 0--2.
		\item \textbf{Cirugía:} la consideración se basa en la madurez esquelética evaluada.
		\item \textbf{Observación:} pacientes con Risser 4--5 requieren seguimiento menos
		intensivo.
	\end{itemize}
	
	\subsubsection*{Correlaciones y Consideraciones}
	
	\textbf{Edad y género:} las niñas alcanzan la madurez esquelética antes que los
	varones. La edad cronológica no siempre correlaciona con la madurez ósea; los factores
	genéticos influyen en la velocidad de osificación.
	
	\textbf{Otros indicadores de madurez:} estadios de Tanner (desarrollo sexual
	secundario), edad ósea de mano por método de Greulich-Pyle, y evaluación del
	cartílago trirrradiado pélvico.
	
	\textbf{Limitaciones del método:}
	\begin{itemize}
		\item Variabilidad entre poblaciones: diferencias étnicas en patrones de osificación.
		\item Progresión rápida: los estadios 1--4 ocurren aproximadamente en un año.
		\item Detección tardía: la osificación puede aparecer hasta 18 meses después del
		pico de crecimiento.
	\end{itemize}
	
	\textbf{Sistemas alternativos:} clasificación de Sanders (radiografías de mano),
	ultrasonido para detección más temprana de osificación, y el sistema Risser modificado
	que incorpora el cartílago trirrradiado.
	
	\textbf{Diferencias regionales:} el sistema norteamericano y el europeo presentan
	variaciones en la definición del estadio Risser 4, lo que puede afectar las decisiones
	de tratamiento.
	
	\begin{figure}[h]
		\centering
		\includegraphics[width=0.7\linewidth]{imagenes/pelvis46}
		\caption{}
		\label{fig:pelvis46}
	\end{figure}
	
	
	%----------------------------------------------------------
	%  Bloque informativo — no duplicar con \subsection de enfoque
	%----------------------------------------------------------